\documentclass[parskip=full]{article}

%%%%%%%%%%%%%%%%%%%%%%%%%%%%%%%%%%%%%%%%%%%%%%%%%
% Basic includes
%%%%%%%%%%%%%%%%%%%%%%%%%%%%%%%%%%%%%%%%%%%%%%%%%


\usepackage[pdftex, colorlinks=true, linkcolor=blue, citecolor=blue, urlcolor=blue, bookmarks=false]{hyperref}

\usepackage{arydshln}
\usepackage{array}
\usepackage{makecell}
\usepackage{tikz}
\usepackage{amsmath}
\usepackage{amsfonts}
\usepackage{verbatim}
\usepackage{mathrsfs}
\usepackage{amssymb}
\usepackage{amsthm}
\usepackage{bbm}
\usepackage{bm}
\usepackage[capposition=top]{floatrow}
%\usepackage[dvipsnames]{xcolor}
\usepackage{ulem}
\usepackage[labelfont=bf]{caption}
\usepackage{subcaption}
\usepackage{graphicx}
\usepackage{float}
\usepackage{natbib}
\usepackage{booktabs}
\usepackage{longtable}
\usepackage{multirow}
\usepackage{rotating}
\usepackage{pdflscape}
\usepackage[title]{appendix}
\usepackage{footmisc}
\usepackage{bigdelim}
\usepackage{tabularx}
\usepackage{xspace}
\newcommand{\A}{\ensuremath{\mathcal{A}}\xspace}
\newcommand{\B}{\ensuremath{\mathcal{B}}\xspace}
\newcommand\pa[1]{\ensuremath{\left(#1\right)}}
\usepackage{epigraph}
\usepackage{setspace}
%\onehalfspacing
\setstretch{1.5}

%%%%%%%%%%%%%%%%%%%%%%%%%%%%%%%%%%%%%%%%%%%%%%%%%%%
% Microtype
%%%%%%%%%%%%%%%%%%%%%%%%%%%%%%%%%%%%%%%%%%%%%%%%%%%
\usepackage[activate={true,nocompatibility},final,tracking=true,kerning=true,spacing=true,factor=1100,stretch=10,shrink=10]{microtype}
% activate={true,nocompatibility} - activate protrusion and expansion
% final - enable microtype; use ``draft'' to disable
%tracking=true, kerning=true, spacing=true - activate these techniques
% factor=1100 - add 10% to the protrusion amount (default is 1000)
% stretch=10, shrink=10 - reduce stretchability/shrinkability (default is 20/20)

%%%%%%%%%%%%%%%%%%%%%%%%%%%%%%%%%%%%%%%%%%%%%%%%%%%
% Hyperref setup
%%%%%%%%%%%%%%%%%%%%%%%%%%%%%%%%%%%%%%%%%%%%%%%%%%%
\usepackage{url}
\makeatletter
\g@addto@macro{\UrlBreaks}{\UrlOrds}
\makeatother
\usepackage{hyperref}
\hypersetup{
  colorlinks,%
  citecolor=black,%
  filecolor=black,%
  linkcolor=blue,%
  urlcolor=black
}

%%%%%%%%%%%%%%%%%%%%%%%%%%%%%%%%%%%%%%%%%%%%%%%%%%%
% Font setup
%%%%%%%%%%%%%%%%%%%%%%%%%%%%%%%%%%%%%%%%%%%%%%%%%%%
\usepackage{mathpazo} % math & rm
%\linespread{1}        % Palatino needs more leading (space between lines)
\usepackage{parskip}
\usepackage[scaled]{helvet} % ss
\usepackage{courier} % tt
\usepackage{natbib}

\usepackage{filecontents}

\normalfont
\usepackage[T1]{fontenc}

\usepackage{titlesec}
\titleformat{\section}
{\normalfont\fontsize{12}{17}\bfseries}
{\thesection}
{1em}
{}

\titleformat{\subsection}
{\normalfont\fontsize{11}{16}\bfseries}
{\thesubsection}
{1em}
{}

%%%%%%%%%%%%%%%%%%%%%%%%%%%%%%%%%%%%%%%%%%%%%%%%%
% macros
%%%%%%%%%%%%%%%%%%%%%%%%%%%%%%%%%%%%%%%%%%%%%%%%%
\newcolumntype{C}[1]{>{\centering\let\newline\\\arraybackslash\hspace{0pt}}m{#1}}
\newcommand{\indicator}[1]{\mathbbm{1}_{\left\{ {#1} \right\} }}
%\newcommand{\resultsloc}{../code/python/output/}
\newcommand{\subfigwidth}{0.475\textwidth}
\newcommand{\subfigwidthtwo}{0.375\textwidth}
\newcommand{\figwidth}{0.6\textwidth}
\DeclareMathOperator*{\argmax}{arg\,max}

\setlength\epigraphwidth{0.9\textwidth}
\setlength\epigraphrule{0pt}
\renewcommand{\epigraphflush}{center}
\usepackage{etoolbox}
\makeatletter
\patchcmd{\epigraph}{\@epitext{#1}}{\itshape\@epitext{#1}}{}{}
\makeatother

%%%%%%%%%%%%%%%%%%%%%%%%%%%%%%%%%%%%%%%%%%%%%%%%%
% Page and paragraph settings
%%%%%%%%%%%%%%%%%%%%%%%%%%%%%%%%%%%%%%%%%%%%%%%%%
\setlength{\parskip}{0.0pc} \setlength{\parindent}{0.25in}
\setlength{\topmargin}{-3pc}
\setlength{\textheight}{9in}
\setlength{\oddsidemargin}{0pc} \setlength{\evensidemargin}{0pc}
\setlength{\textwidth}{6.5in}
\makeatletter
\@fpsep\textheight
\makeatother

\newenvironment{fignote}{\begin{quote}\footnotesize}{\end{quote}}
\def\sym#1{\ifmmode^{#1}\else\(^{#1}\)\fi}
\renewcommand{\baselinestretch}{1.5}

% Formats the notes at the bottom of tables and figures.
\newenvironment{fignote2}[1]
{
\par
\medskip
\begin{minipage}{#1}
\footnotesize
}
{
\end{minipage}
}

% Formats the notes at the bottom of tables and figures.
\newenvironment{fignote3}[1]
{
%\par
\begin{minipage}{#1}
\footnotesize
}
{
\end{minipage}
}

%%%%%%%%%%%%%%%%%%%%%%%%%%%%%%%%%%%%%%%%%%%%%%%%%
% Title page
%%%%%%%%%%%%%%%%%%%%%%%%%%%%%%%%%%%%%%%%%%%%%%%%%
%%%%%%%%%%%%%%%%%%%%%%%%%%%%%%%%%%%%%%%%%%%%%%%%%
% Title page
%%%%%%%%%%%%%%%%%%%%%%%%%%%%%%%%%%%%%%%%%%%%%%%%%

\title{\vspace{-1.8cm}The Welfare Consequences of Rent Controls\footnote{We extend our gratitude to the seminar attendees at the University of Wisconsin-Madison, University of Southern California, Queen's University, York University, McMaster University, University of Windsor, as well as to Efraim Benmelech, Boaz Abramson, Pau Pujalos, Shouyong Shi, Edward Glaeser, Joseph Steinberg, Amy Sun, Yahong Zhang, Ruoyu Chen, Chenggang Zhou, and John Baker for their valuable insights and comments. We owe a special thanks to Joseph Steinberg for generously providing us access to his high-powered computer for this research. The data and replication code are available on GitHub at \url{www.github.com/srotberg/the_welfare_consequences_of_rent_controls}. For those who like audio, we have an AI-generated Spotify episode explaining the paper: \url{https://spotifycreators-web.app.link/e/QVkPeRFoUOb}.}}

\author{Shahar Rotberg\footnote{Canada Mortgage and Housing Corporation. Corresponding author contact: \url{shahar.rotberg@gmail.com}. Website: \url{www.sites.google.com/site/shaharrotberg}. Shahar Rotberg did this work on his personal equipment and time. He exclusively owns the code that was used to solve the model in this paper.} \and Lin Zhang\footnote{Canada Mortgage and Housing Corporation. Contact: \url{lzhang@cmhc-schl.gc.ca}. The opinions expressed herein do not necessarily represent the views of CMHC.}}


\pagestyle{plain}

\begin{document}

\maketitle

\setstretch{1.5}
\begin{abstract}
\onehalfspacing
\noindent

\noindent Unit control and tenancy control are prevalent rent regulations that limit rent hikes. Unit control offers immediate discounts but creates shortages, resulting in controlled units being allocated via lotteries. Tenancy control allows landlords to adjust rents upon tenant turnover, but generates no shortages, because a price mechanism allocates controlled units. Integrating both policies into a quantitative general equilibrium model of housing calibrated to Toronto, we find that tenancy control yields greater welfare than unit control and reduces misallocation of controlled units. Moreover, the maximizing-welfare degree of tenancy control is 65\%. 
%words=90
\\[1.5ex]

\noindent \textbf{JEL Classification:} E17; R31; R52. 

\noindent \textbf{Keywords:} General equilibrium, housing, rent control, tenancy control, unit control, homelessness.
\end{abstract}
\setstretch{1.5}

\section{Introduction}

Tenancy control and unit control are the two most common rent control policies implemented worldwide to keep rental housing affordable. Unit control, seen in cities like New York City, Vancouver, and Paris, restricts rent increases for rent-controlled units both during a tenancy and when tenants change. In contrast, tenancy control, applied in cities such as Toronto, San Francisco, and Hong Kong, regulates rent increases only during a tenancy, but allows landlords to set new rents freely when tenants vacate. While the economic literature has extensively examined unit control, it has left tenancy control largely under-explored, especially in general equilibrium. 

This paper proposes the first general-equilibrium theory of tenancy control and delves into essential distinctions between these two types of rent control. We construct a quantitative general-equilibrium housing model that nests both tenancy control and unit control, and using the model, we evaluate the welfare implications of both of these rent control policies. Our quantitative exercises using data from Toronto show that moving from tenancy control to unit control would reduce welfare and worsen the allocation of controlled units. Furthermore, welfare is maximized when 65\% of the rental market is subject to tenancy control.

Unit control and tenancy control differ in the way they allocate controlled units. Unit control offers new tenants immediate access to discounted rents compared to free-market rents. However, this creates a shortage of rental housing because the quantity demanded surpasses the quantity supplied at the government-mandated rent ceiling. Consequently, as \citet{gl-2003} argue, rental units are assigned arbitrarily, and thus do not necessarily go to households with the highest willingness to pay. Conversely, tenancy control permits landlords to adjust rents when tenants change, and thus, when there is no minimum restriction on rental housing consumption, it ensures that market forces allocate controlled units to households with the highest willingness to pay. This mechanism averts shortages, but it entails that new tenants must pay a premium over free-market rents to secure controlled accommodation, thereby postponing the advantages of rent control for tenants.

The ability of the price mechanism under tenancy control to improve the allocation of rental housing is brought into question, however, because in the real world housing cannot be consumed in arbitrarily small quantities. Thus, when renters are required to pay a premium to secure a controlled unit under tenancy control, certain low-income households with the highest willingness to pay (highest marginal utilities) may be unable to afford access the controlled market. Therefore, whether tenancy control results in a better allocation of controlled units relative to unit control is a quantitative question.

We construct a model that encompasses both types of rent-control policies, with tenancy control serving as our baseline model. Our model is a general-equilibrium housing framework featuring overlapping generations and two rental markets operating in parallel: an uncontrolled market (i.e., free market) and a controlled market. Within our baseline model, which features tenancy control, renters have the freedom to choose between renting in the controlled market or the free market. In the controlled market, rents for units are regulated throughout a renter's tenancy, gradually declining in real terms as the renter remains longer in the same unit. However, these units also deteriorate due to landlords' lack of maintenance. Upon a tenant's departure, landlords can revert to charging market rents for new tenants. Conversely, units in the free market are unregulated. Landlords maintain these units adequately and rents do not decrease in real terms regardless of tenancy duration. Moreover, rents in this market are subject to stochastic fluctuations, which renters are unable to insure themselves against.

To capture the barriers hindering certain low-income households from entering into rental markets, we introduce a minimum rental size, and account for the possibility of households facing homelessness or residing with their parents.\footnote{Homelessness is critical to our model because it matters for how efficiently controlled units are allocated. Our model also includes the option for young households to live rent-free with their parents for three reasons. Firstly, and most importantly, this element determines the welfare impact of barriers to entry in the controlled market. Ignoring this aspect would lead us to falsely categorize many low-income but non-rent-burdened young households as rent-burdened or homeless, and thus distort our welfare results. Secondly, this option influences young households' decisions regarding renting and owning versus living with their parents. In its absence, some housing market outcomes, such as the impact of full decontrol on free market rents, in our counterfactual analyses are qualitatively opposite. Thirdly, almost a quarter of 26-34 year-olds live with their parents, which is a substantial portion that, as we explain above, affects housing market outcomes.}\footnote{Although we do not have data on this for Toronto, according to estimates by \citet{spatial_distance-2020}, approximately 75\% of young adults (aged 26-34) in the U.S. reside within a 30-mile radius (approximately 48 kilometers) of at least one of their parents. This suggests that similar numbers most likely apply to Toronto.} Building upon \citet{az-2024}, households experiencing homelessness do not pay rent but receive minimal housing services, akin to shelter assistance local governments typically provide. Similarly, households cohabiting with their parents reside rent-free, albeit with reduced housing consumption compared to what they could secure in the rental market.

Our framework yields three primary equilibrium implications in our baseline model with tenancy control. Firstly, to capitalize on the gradual reduction in rents over time and the insurance against rent fluctuations offered by rent control, tenants must initially pay a premium to rent in the controlled market. Secondly, when the government reduces the size of the controlled market, intensified competition for rent-controlled units pushes the premium upward. Thirdly, high-income households, equipped to handle the rent fluctuations in the free rental market and anticipating to become homeowners quickly, often opt for the free market to avoid paying the premium. Conversely, low-income households, with lower risk tolerance, an expectation of prolonged rental tenure, and the ability to enter the controlled market, typically prefer to pay the premium to enjoy the advantages of rent control.

The remainder of our model encompasses fairly standard components. Households make decisions regarding renting, homeownership, or, as previously discussed, homelessness or residing with their parents. Homeowners can access long-term mortgage contracts. Working-age households earn labor income by supplying labor inelastically and receive pensions upon retirement. Households accrue capital income from their financial assets, and labor and capital incomes are subject to progressive taxation. Similar to the framework outlined in \citet{garrigaeetal-2009}, a price-taker rental company, operating with convex marginal costs, supplies rental housing. Drawing inspiration from \citet{fallis-1985}, we model rental housing supply with exemptions after the structure observed in Toronto's rental market, where older construction is subject to rent control while newer construction is not. In this configuration, the average free-market rent represents the rental company's marginal revenue, thereby determining the quantity of rental housing supplied for profit maximization. Additionally, mirroring the approach outlined in \citet{rotberg-2022}, a price-taker housing construction sector, characterized by convex marginal costs, contributes to the creation of new housing units, which are subsequently added to the depreciated housing stock.\footnote{Households save by investing in bonds and stocks of rental and construction firms. Our economy is a city with an exogenous risk-free interest rate. Thus, a no-arbitrage condition applies, which means that households are indifferent between these options. Since we only consider long-run equilibria, changes in the stock prices of the rental and construction firms between equilibria do not affect our results.}

The model incorporates both federal and local government entities. The federal government is responsible for collecting consumption and income taxes, and utilizes the generated revenue to fund pensions for retired households. Conversely, the local government collects property taxes and allocates the proceeds towards providing homeless shelters, as in \citet{abramson-2024}, and a public good that does not directly impact household utility. Given that rent control typically falls under the jurisdiction of local governments, we require budget balance only at the level of the local government, which we achieve by adjusting the property tax rate.

We calibrate the long-run equilibrium of our model using data from Toronto, Canada, where, in 2016, a portion of the rental market was subject to tenancy control while the remainder of the rental market remained free. Much of our calibration process follows standard practices; we select parameter values to align with various statistics such as the homeownership rate, and the proportion of total renters' income allocated to rent, among others. However, certain aspects of our calibration are more unique to our study. Firstly, we target the percentage of renters in the controlled market, which stands at 71\% according to the 2016 Census. This figure serves as a crucial determinant for the premium that new tenants in the controlled market must pay to secure a unit in equilibrium. Secondly, we target the percentage of renters spending over 50\% of their income on rent, which we also draw from the 2016 Census. This target, akin to the approach proposed by \citet{rs-2023}, informs the determination of the minimum rental unit size, and thus reflects the barrier faced by low-income households in entering the rental market. Lastly, we set the level of housing services obtained by households living with their parents and those experiencing homelessness to correspond with the percentage of 26-34-year-old households cohabiting with parents and the percentage of homeless households, respectively. We derive the former from the 2016 Census, while we source the latter from '`The Homelessness Hub''.\footnote{See \url{www.homelesshub.ca}.}

In our exercises, our general approach is to contrast our calibrated baseline model featuring two rental markets operating in parallel (our benchmark) with alternative long-run equilibrium scenarios, which we expand on in the paragraph below. Drawing from the principles outlined in \citet{ralws-1971}, we define welfare as the ex-ante expected utility of newborns, a widely accepted convention in the literature \citep{conesa-et-al,gkkoc}.

In our initial experiment, we juxtapose our benchmark model with tenancy control against unit control. Under unit control in our model, households enter a lottery for an opportunity to rent in the controlled market. Those with an opportunity and who opt for the controlled market can retain their unit indefinitely. Conversely, households lacking this opportunity may choose to rent in the free market, pursue homeownership, reside with their parents if they can, or face homelessness. To ensure comparability with our baseline, we implement two adjustments. Firstly, we establish the controlled rent in the unit control scenario, mandated exogenously by the government, to mirror the average controlled rent in the baseline's benchmark. Secondly, we equate the proportion of the rental market under control in the unit control setup to the proportion of the rental market under tenancy control in our benchmark.\footnote{We also align the rent reduction in the unit control setup with the rate of unit deterioration resulting from landlord neglect. This approach ensures that controlled units deteriorate at an equivalent pace across both scenarios while maintaining constant real controlled rent in our unit control experiment.}

Overall, moving from tenancy control to unit control would raise aggregate rental demand, average free-market rents, and homelessness, and slightly reduce welfare. The reasons for these results are as follows. Because unit control would result in a shortage of controlled units, these units would be allocated using lotteries. We find that quantitatively, despite the minimum rental size, fewer controlled units would be allocated to low-income households compared to our benchmark with tenancy control and thus the homelessness rate would rise substantially. However, new tenants in the controlled market would face substantial rent discounts compared to the benchmark, thus increasing aggregate rental demand and prompting a 2.4\% escalation in average free-market rents and a slight increase in house prices of 0.43\%. Despite the worse misallocation of controlled units, and higher average free-market rents and house prices, newborns' welfare would only decline by 0.06\%. This is because unit control would offer immediate rent discounts that take approximately 11 years for a tenant under tenancy control to attain. 

On the supply side, the increase in house prices would raise the cost to supply rental housing, and thus rental supply would shift leftward. However, the increase in average free-market rents would lead to a movement up along the rental supply curve that would outweigh the leftward shift, leading to an expansion in the quantity of rental housing supplied. We offer additional intuition using Figure \ref{fig:diagram_unit_control}.

Lastly, we also demonstrate that the welfare gains associated with the immediate rent discounts under unit control would outweigh the welfare losses from the worse misallocation of controlled units \citep{gl-2003}. Welfare would actually slightly decline under unit control due to the increase in average free-market rents.

Our second experiment compares the baseline benchmark with the long-run equilibrium wherein the government fully decontrols the rental market. Our main results are that full decontrol would raise average free-market rents, aggregate rental demand, house prices, and homelessness, and reduce welfare. The explanation for these effects is the following. Full decontrol would precipitate two primary effects. Firstly, it would diminish the appeal of renting, thereby increasing homeownership, driving up house prices, and reducing aggregate rental demand --- a consequence stemming from the extensive margin. Conversely, full decontrol would reduce the initial rent for new tenants previously renting in the controlled market because average free-market rents are lower than initial controlled rents due to the premium new tenants must pay to access the controlled market. This lower initial rent would spur an increase in aggregate rental demand from those tenants remaining in the rental market --- a repercussion arising from the intensive margin. Overall, the impact from the intensive margin would marginally outweigh that from the extensive margin, resulting in a slight long-run rise in average free-market rents. 

Moreover, the removal of the controlled market as an option for low-income households would also expose more of them to higher average rents and rent volatility, thus raising the homelessness rate even more significantly than a switch to unit control would. The combination of eliminating the controlled market option for certain renters, the marginally higher average free-market rents, and elevated house prices would collectively lower welfare. 

On the supply side, the increase in average free-market rents would exactly offset the leftward shift in rental supply stemming from the higher house prices, leading to no change in the equilibrium quantity of rental housing supplied. We provide further visual intuition in Figure \ref{fig:diagram_full_decontrol}.

Our third experiment contrasts the benchmark with the long-run equilibrium scenario in which all rental units are subject to tenancy control. On the whole, full tenancy control would raise average controlled rents but reduce average rents, house prices, aggregate rental demand, homelessness, and welfare. These results would arise because of the following reasons. First, when the entire market is under tenancy control, no households would be exposed to free-market rent fluctuations, and thus the homelessness rate would marginally decline.

Second, extending tenancy control across the entire market would also trigger extensive and intensive margin effects, albeit in reverse compared to the second experiment. At the extensive margin, an increased number of high-income households would opt for renting, consequently reducing homeownership rates and house prices while increasing aggregate rental demand. Conversely, at the intensive margin, the higher initial rent imposed on renters who previously leased in the free market would diminish their demand for rental housing. Although the effect at the intensive margin would once again outweigh that at the extensive margin, aggregate rental demand for controlled units would exceed that in the benchmark. 

Because more high-income households would rent in the controlled market, and these households tend to transition to homeownership more quickly than households who rented in the controlled market in the benchmark, the average controlled rent would rise, but initial rents would remain relatively unchanged. Overall, it is the removal of the option to rent in the free market that would lead to small welfare losses. 

In contrast to the previous two experiments, here the decrease in house prices would lower the cost of providing rental housing, and thus shift rental supply to the right. However, the decrease in the marginal rent, which would now be the average rent in the controlled market, would cause a movement down along the rental supply curve that would outweigh the rightward shift in rental supply, and thus result in a decrease in the quantity of rental housing supplied in equilibrium. Figure \ref{fig:diagram_full_control} provides further visual intuition.

In our fourth experiment, we endeavor to identify the degree of tenancy control that maximizes welfare in our baseline model. We vary the proportion of rental housing subjected to tenancy control and assess welfare in the long-run equilibrium of each scenario. Our analysis reveals that, in our model, the welfare-maximizing degree of tenancy control stands at 65\%, a figure remarkably close to our benchmark where approximately 70\% of the rental market falls under tenancy control. Low degrees of tenancy control prove detrimental to welfare as they exacerbate the scarcity of controlled units, driving up their premiums and further excluding low-income households from the market. Conversely, high degrees of tenancy control also reduce welfare. As the share of the controlled market increases past a certain threshold, average free-market rents need to rise to incentivize more households to opt for rentals within the controlled sector. However, this dynamic adversely affects households that favor renting outside the controlled sector.\footnote{Consistent with the findings of our first experiment, we ascertain that unit control would diminish welfare under the same degree of control and the same average controlled rents as in our welfare-maximizing long-run equilibrium.}

This paper makes significant theoretical, methodological, and quantitative contributions to the field of rent control research. Theoretically, we uncover fundamental economic distinctions between tenancy control and unit control, which, to our knowledge, have not been previously identified or discussed in the literature. Methodologically, we develop a general equilibrium model that includes both tenancy control and unit control, and thus provide a versatile framework for analysis. Quantitatively, we leverage our model to assess the welfare implications of both types of rent controls, shedding new light on their comparative effectiveness. Furthermore, we pinpoint the welfare-maximizing degree of tenancy control in our model, and thus offer valuable insights into policy optimization in this domain.

\section{Related literature}\label{sec:related_lit}

The classical theory in economics, discussed in \citet{suen-1989}, posits that government-imposed rent ceilings (termed unit control in our study) cause demand to surpass supply, and thus induce a shortage of controlled units. \citet{gl-2003} contend that when such shortages arise, the allocation of controlled units is done through lotteries, which lead to a misallocation. However, we argue in our paper that this classical theory does not hold true for tenancy control. With tenancy control, also known as vacancy decontrol or the third generation rent control, landlords can adjust rents freely for new tenants. Consequently, competition among potential tenants for these units raises their price relative to free-market units. This means that no shortages of controlled rental housing exist. Moreover, so long as the combination of the rent premium and minimum rental size exclude fewer low-income households from accessing controlled units relative to unit control, tenancy control would lead to a better allocation of controlled units. 

As highlighted in \citet{mm-2023}, the scarcity of empirical research on rent control arises from two main factors: ``First, micro-level data on house prices and rents is mostly available for recent periods, whereas sources of potentially exogenous variation are available for earlier episodes. Second, there are various ways in which rental cap policies are implemented, and hence, we have only a partial understanding on how certain types of rental cap policies may affect the market." Therefore, it is imperative to employ equilibrium housing models to learn about the impacts of different rent-control policies. \citet{fallis-1985} examined Toronto's rent-control policies, albeit using a static model rather than a dynamic one like ours. Similar to our model, theirs included an interaction between controlled and free markets. More recent studies have turned to general equilibrium models with dynamic lifecycle choices. \citet{eric-1995} and \citet{fmn-2023} model rent control akin to unit control and assess its effects on house prices, rents, and welfare.\footnote{Controlled units in \citet{fmn-2023} offer immediate discounts that increase with tenure duration.} However, neither of them incorporates tenancy control as we do, nor do they compare tenancy control to unit control. Moreover, although \citet{ms-1993} show that rent control leads to a deterioration in the quality of rental housing, neither \citet{eric-1995} nor \citet{fmn-2023} allow controlled units to deteriorate as tenants remain longer in their unit.

One element that is included in \citet{fmn-2023} but is omitted in our analysis is the consideration of migration between cities with unit control and those without. They model New York City as a city with unit control, contrasting it with neighboring areas lacking such regulations altogether. However, since Ontario's tenancy control laws apply uniformly at the provincial level, encompassing all neighboring cities of Toronto, this spatial element is not applicable to our investigation.

\begin{comment}

Save for later if needed. 
Furthermore, as highlighted by \citet{statcan-2024}, working couples in Canada's top 10% income bracket are nine times more likely to telecommute compared to those in the bottom 10% income bracket. Additionally, over 50% of jobs in Toronto are deemed suitable for remote work. Given these statistics, the influence of spatial factors on our findings may be quite negligible.

\end{comment}

Our work is also related to a couple of theoretical studies with a focus on modeling tenancy control. \citet{be-2000} model tenancy control in an asymmetric information setup in which landlords do not know the types of tenants or how long tenants stay. While their model is focused on strategic interactions between landlords and tenants, our work provides a theory on how controlled units are allocated when free and tenancy-controlled markets operate at the same time, and we compare it to unit control. In addition, \citet{diamond-2019} evaluate the welfare impact of expansions in tenancy control in San Francisco using a dynamic, neighborhood choice model. Their model features many desirable features, such as tenant buyouts, the potential for biased renters' expectations about the future, and endogenous landlord responses. However, it does not feature general equilibrium elements that we focus on, such as endogenous rents, feedback effects between free-market rents and controlled rents, or interactions between the rental and homeownership markets. Furthermore, our paper introduces a novel theoretical perspective that diverges from the classical economic theory of rent control by demonstrating that tenancy control does not result in shortages.

Our theory provides a theoretical explanation for the rent premiums in controlled markets observed by several studies. \citet{nagy-1997} documented that rent-controlled apartments initially commanded higher rents than free-market ones, while \citet{early-2000} and \citet{cjq-2023} estimated that tenants in the controlled market were more likely to pay rent premiums during their initial years of tenancy. Although all these studies examined rent control in New York City where unit control is in place, it suggests that in many cases, landlords may be breaking the law and actually resetting rents upon tenant turnover. When viewed from the perspective that, in practice, unit control may actually result in equilibria that are closer to tenancy control \citep{be-2000}, our model simulations of full decontrol may also align with findings from empirical studies on Cambridge, Massachusetts where unit control was repealed. For instance, \citet{sims-2007} observed that the decontrolling in Cambridge in 1995 led to rent increase and \citet{app-2014} found that house prices also increased following this decontrol. 

Our results from the experiment on full tenancy control relate to some findings from \citet{dmq-2019}, who use a structural model to evaluate the impacts of the expansion of tenancy control in San Francisco in 1994 that nearly covered the entire rental market and only exempted new construction of rental buildings with 4 units or fewer. This suggests, that in practice, the marginal rent in San Francisco following this expansion was likely the average controlled rent, since multi-unit buildings (5 units and more), which constituted the majority of rental housing were not exempted. Their model estimates suggest that this expansion of tenancy control reduced the quantity of rental housing supplied, raised average controlled rents, and reduced young household's welfare. When we move from our model's benchmark to full tenancy control, average controlled rents rise, the quantity of rental housing supplied declines, and newborns experience a welfare loss. 

Lastly, our work also contributes to the literature on the relationship between rent control and homelessness. First, our simulation results align with the empirical findings of \citet{eo-1998}, who estimated that more rent control reduces homelessness. Second, while \citet{az-2024} and \citet{abramson-2024} have developed general equilibrium models of homelessness, to the best of our knowledge, we are the first to integrate homelessness into a general equilibrium model with tenancy and unit controls, and investigate the connection between these two types of rent control and homelessness.

\section{Institutional background}\label{sec:insitutional_background}

Similar to the United States and certain European nations, Canada also introduced rent controls as a wartime measure. During World War II, the National Housing Act of 1944 introduced rent controls in Canada. This initial wave of rent control, sometimes referred to as the ``first-generation", entailed freezing nominal rents and was inspired by similar measures enacted in European countries during World War I. The primary objective was to shield veterans from evictions and rent hikes during wartime pressures.

These rent controls seemed as a pragmatic public policy response during a period of heightened demand for rental units, coupled with restrictions on new construction to conserve materials for war efforts \citep{downs-1988}.\footnote{In a similar vein, rent freezes were imposed temporarily during the peak of the COVID-19 pandemic.}

While the rent freeze appeared warranted during wartime, many economists voiced concerns about its lasting impacts. These concerns included diminished rental supply, reduced labor mobility, and compromised property maintenance \citep{arnott-1995}. Consequently, the Ontario provincial government repealed the ``first-generation" of rent controls in 1949.

In the 1970s, rent controls resurfaced in response to rapid inflation, echoing actions taken by other jurisdictions. In 1975, against the backdrop of the oil crisis, the Ontario government reintroduced rent controls through the Residential Premises Rent Review Act. These renewed controls, characterized as the ``second-generation" rent controls or unit controls, diverged from their ``first-generation" predecessors in three key ways.

Firstly, they permitted modest nominal rent increases annually, guided by prescribed directives. Secondly, they incorporated cost pass-through clauses, enabling courts to approve increases beyond statutory guidelines in cases of significant rises in the costs associated with supplying rental housing, landlord financial losses, or necessary capital investments. Lastly, these controls exempted luxury accommodations, government-owned buildings, social housing, and private rental units constructed after January 1976.

Since 1975, the Ontario rent control system has undergone significant evolution, marked by two major changes. The first notable shift occurred with intermittent exemptions. In 1985, the Residential Rent Regulation Act eliminated exemptions for previously exempted units and newly constructed units. This change subjected both new constructions and buildings erected after 1976 to rent controls. However, in 1992, the Rent Control Act reintroduced rent control with exemptions for new constructions. Under this act, all rental buildings erected after November 1, 1991, were exempt from rent controls. Nonetheless, in April 2017, the Ontario government reversed this exemption policy by extending rent controls to all rental units constructed after November 1, 1991. Subsequently, in 2018, the Ontario government introduced new legislation that reinstated exemptions, but with a revised scope. Under this legislation, only buildings constructed or first occupied after November 15, 2018, were exempt from rent controls.

The second significant transformation in the Ontario rent control system occurred in 1998 with the implementation of tenancy control under the Ontario Tenant Protection Act. This legislation granted landlords of rent-controlled units the ability to increase rents beyond the limits prescribed by rent increase guidelines, but exclusively when entering into a new lease agreement with a new tenant. Otherwise, landlords were mandated to adhere to the stipulated rent increase guidelines. This form of rent control, often referred to as the ``third-generation" rent control, introduced a nuanced approach to regulating rental prices \citep{arnott-2003}. Since its enactment, the Ontario government has not implemented any further changes in this regard.

Toronto's experience with rent controls reflects the evolution observed in many other cities, spanning the three generations of regulatory frameworks. It transitioned from rent freezes during wartime to decontrol in 1949, followed by unit control in 1975, and eventually to tenancy control in 1998. Despite the diversity of rent control models worldwide nowadays, rent controls today can generally be classified into two main categories: stringent rent controls that correspond to unit control and temperate rent controls that correspond to tenancy control. However, both these categories also share common features such as exemptions, cost pass-through clauses, and restrictions on condominium conversions. 

For example, rent control laws in New York City fall under the umbrella of unit controls, featuring exemptions, cost pass-through clauses, and restrictions on condominium conversions. By employing this classification method, rent-control policies in other cities like Vancouver, Montréal, and Paris can also be categorized as unit controls. Conversely, the category of tenancy controls includes cities such as San Francisco, Los Angeles, Washington D.C., Singapore, and Hong Kong, to list some. These jurisdictions implement rent-control policies that focus on regulating rental prices during a tenant's occupancy, with varying exemptions and regulations tailored to each locality's unique housing market dynamics.

Amidst rising housing costs and overall inflationary pressures, numerous jurisdictions worldwide have responded by enacting or revising rent control legislation. In Canada, Nova Scotia implemented rent controls in 2020, joining a growing list of provinces addressing housing affordability concerns. Across the United States, states including Oregon, California, New Jersey, Maryland, and the District of Columbia introduced rent controls in either 2019 or 2020, reflecting a trend towards regulatory intervention in housing markets. Similarly, in Europe, cities such as Paris and Berlin, and the Catalonia region saw the (re)introduction of rent controls during the same period, highlighting a concerted effort to address affordability challenges in urban centers. Even in developing countries, the adoption of rent control measures has gained traction, with China implementing such regulations for the first time in 2021. These developments underscore the widespread recognition of the need to address housing affordability and tenant protections across diverse global contexts, but also the need to model rent controls and evaluate their impacts on household well-being.

\section{Model}\label{sec:model}

The model encompasses overlapping generations of finitely-lived households, price-taker rental and housing construction firms, and governmental bodies. Each age cohort of households exhibits heterogeneity in labor market productivity, housing, mortgage debt, financial assets, and housing tenure (ownership, renting, cohabiting with parents, or homelessness). Renters further differ in rental market type (controlled\footnote{A couple of factors may explain why renters, despite the potential benefits, rarely engage in contracts to regulate their rents when markets are unregulated. First, the upfront legal costs of establishing and enforcing individual contracts can be prohibitive, especially for low-income households who would benefit from rent control the most. Second, landlords might hold greater bargaining power than renters, especially in expensive cities, making it difficult for renters to negotiate favorable terms. In this context, government regulation of rental markets is a mechanism that eliminates the upfront costs of contracting and enhances renters' bargaining power.} or free), and the duration of their residence in a particular unit. Renters in the free market face rent shocks.

Pre-retirement, households engage in the labor market, while post-retirement, they receive government pensions dependent on their average lifetime labor income. Renters make decisions regarding rental market selection and tenure duration. Homeowners incur proportional property taxes and maintenance expenses, and carry mortgages as long-term contracts, following the approach of \citet{kkk-2021}. Until age $J_{p}$, households may opt for rent-free accommodation with their parents, albeit at the cost of reduced housing services compared to renting in the rental market. Past age $J_{p}$, households cannot cohabit with their parents, and thus face homelessness if they are unable or unwilling to engage in the rental market, receiving minimal housing services at no cost.

Competitive rental firms, characterized by convex marginal costs, operate in both the controlled and free rental markets simultaneously. Similarly, competitive construction companies, subject to increasing marginal costs, invest in new housing production, which supplements the depreciated housing stock.

The model integrates both federal and local government entities. The federal government finances pensions for retirees through income and consumption taxes, whereas local governments offer public services to residents and shelter services to homeless households, which they fund by imposing property taxes. Given the localized nature of rent control legislation, we require budget balance solely at the local government level.

Considering that our model economy is a city, we assume that the interest rate, $r$, is exogenous \citep{sommeretal-2013, arslanetal-2015, sommersullivan-2017, kaplanetal-2019}. Household savings encompass bonds and stock investments in rental and construction firms, with an implicit requirement of indifference between these options due to a no-arbitrage condition. Because asset market clearance is not required, we do not explicitly track dividends associated with stock holdings. Moreover, since we only consider long-run equilibria, changes in rental and construction companies' stock prices have no affect on our results.

\subsection{Demographics}\label{subsection:demographics}

Households enter the model at age $j=1$, retire from the labor force at age $J_{R}$, and live for $J$ periods. Until age $J$, they face no mortality risk, but die with certainty afterward. Household preferences are characterized by the following utility function:

\begin{equation*}
    u_{j}(c,h(o,f,\iota))=\xi_{j}^{\sigma}\frac{\left(c^{1-\gamma} h(o,f,\iota)^{\gamma}\right)^{1-\sigma}}{1-\sigma},
\end{equation*}

\noindent where $c$ denotes non-housing consumption and $h(o,f,\iota)$ represents housing services consumption. This function, which we elaborate on in Section \ref{subsection:housing_choice}, depends on $o$, $f$, and $\iota$. $f=1$ indicates households renting in the free market, while $f=0$ denotes those renting in the controlled market.\footnote{For homeowners, households living with their parents, or homeless households, $f$ holds no significance, and thus we set it to $1$.} The parameter $\iota$ signifies the number of years a renter in the controlled market has resided in the same unit, whereas $o=1$ signifies homeowners and $o=0$ signifies renters, households cohabiting with their parents, or homeless households. The weight $\gamma$ denotes the utility assigned to housing service consumption, $\sigma$ represents the household's constant relative risk aversion, and $\xi_{j}$ scales equivalently to accommodate variations in household size over its lifecycle, drawing from the works of \citet{lm-1980}, \citet{wz-2010}, \citet{xu-2013}, and \citet{rs-2023}. Household discount of future consumption at a rate of $\beta$ and non-housing consumption incurs a proportional tax $\tau_{c}$.

Similar to \citet{cagetti-denardi-2006}, households derive a warm-glow from bequests they leave behind upon death, which we capture by $\omega_{B}(q)=\mu \frac{q^{1-\sigma}}{1-\sigma}$. Here, $q$ represents the household's net worth at the end of their life, which is further defined below, while $\mu$ modulates the strength of the bequest motive.

\subsection{Heterogeneity}

Households are heterogeneous across multiple dimensions:

\begin{itemize}
    \item Financial assets, $a\in \mathcal{A}=\mathbb{R}_{+}$.
    \item Housing tenure, $o \in \mathcal{O}=\{0,1\}$.
    \item Rental market selection for those receiving rental services, $f \in \mathcal{F}=\{0,1\}$.
    \item Duration of residence in the same unit for renters in the controlled market, $\iota \in \mathcal{I}=\mathbb{N}$.
    \item Rent shock in the free market for those renting in the free market, $z \in \mathbb{Z}$.
    \item House size, $h\in \mathcal{H}=\mathbb{R}_{++}$.
    \item Mortgage debt, $m \in \mathcal{M}=\mathbb{R}_{+}$.
    \item Idiosyncratic labor fixed effect $\chi \in \mathcal{X}=\mathbb{R}_{++}$.
    \item Stochastic labor market productivity shock, $e \in \mathcal{E} = \mathbb{R}_{++}$.
\end{itemize}
 
\noindent We denote the measure of age-$j$ households over endowments by $\Psi_{j}(a,o,f,\iota,z,h,m,\chi,e)$. For brevity, we mainly employ $s$ to represent a household's state $(a,o,f,\iota,z,h,m,\chi,e)$.

\subsection{Income and income taxes}\label{subsection:income}

Working-age households ($j<J_{R}$) supply labor inelastically. Age-$j$ households, characterized by labor market productivities $\chi$ and $e$, earn labor income $y_{\ell,j}(\chi,e)=\chi e\zeta_j$. The fixed-effect $\chi$ remains constant throughout a household's life and is drawn from an unconditional distribution, which we denote as $\hat{F}(\chi)$. The idiosyncratic shock $e$ follows a Markov process conditional on the homelessness status, which we denote as $F(e,e'|b)$, where $b=1$ indicates homelessness and $b=0$ indicates non-homelessness. Particularly, in line with \citet{az-2024}, households experiencing homelessness are more prone to encountering persistently low stochastic productivity shocks. $\zeta_j$ represents the deterministic part of households' life-cycle productivity profiles. Retired households receive pensions, $y_{R}(\chi,e)$, which are contingent upon labor productivity at retirement.\footnote{Retirees with higher productivity levels are more likely to have had higher productivities over their lifetimes, and thus receive larger pensions.}

Households can also earn income from their savings, $y_{k}(a)=ar$, where $r$ is the interest rate. We define taxable income as:

\begin{equation*}\label{eq:taxable_labor_income}
  y^{x}_{j}(\chi,e,a) =y_{\ell,j}(\chi,e)+\nu y_{k}(a),
\end{equation*}

\noindent where $\nu$ signifies the taxable portion of capital income. We apply progressive taxes to taxable income as follows:

\begin{equation}\label{eq:income_tax}
  y^{\tau}_{j}(\chi,e,a) = y^{x}_{j}(\chi,e,a)-\tau_{\ell}y^{x}_{j}(\chi,e,a)^{\psi},
\end{equation}

\noindent where $\tau_{\ell}$ determines the average tax rate and $\psi$ governs the degree of progressivity \citep{hsv-2020}.

As in \citet{kkk-2021}, each period, surviving households receive financial wealth endowments equivalent to the bequests left by deceased households in the preceding period. Surviving households across all ages receive financial endowments, which we denote as $\vartheta (\chi,e)$, capturing the correlation between bequests and income.

\subsection{Housing}\label{subsection:housing_choice}

Households have the option to either rent or own a dwelling, denoted by $o \in \mathcal{O}=\{0,1\}$, cohabit with their parents $\left(v \in \mathcal{V}=\{0,1\}\right)$ if $j<J_{p}$, or face homelessness $\left(b \in \mathcal{B}=\{0,1\}\right)$ if $j \geq J_{p}$. Households that opt to rent or own, can also select the size of their dwelling, which we denote by $h$. The smallest rental unit is smaller than the smallest housing unit for purchase, defined as $0<\underline{h}_{r}<\underline{h}_{o}$ \citep{gervais,sommeretal-2013,sommersullivan-2017,kaplanetal-2019}. Households residing with their parents receive housing services equivalent to $\underline{h}_{p}$, while those facing homelessness receive housing services equivalent to $\underline{h}_{h}$, where $0<\underline{h}_{h}<\underline{h}_{p}<\underline{h}_{r}$. Households living with their parents ($v=1$) or experiencing homelessness ($b=1$), receive their housing services free of charge. However, as in \citet{abramson-2024}, the city incurs a cost of $\theta_{b}$ to support each homeless household.

Households opting to own a house of size $h$ incur the purchase price $p h$. Buying and selling houses involve one-time proportional transaction costs, $\tau_{b} p h$ and $\tau_{s} p h$, respectively. In each period, all housing units are subject to property taxation at a rate $\tau_{p}$ and depreciate in value at a rate $\delta_h$.

Renters need to choose between renting in the controlled market or the free market. $f=0$ denotes renting in the controlled market, whereas $f=1$ signifies renting in the free market.

In the free market, renters leasing housing services $h$ pay $p_{r,1} h$. Free-market rent remains unaffected by tenancy duration and is subject to uncertain shocks, $z$, are drawn from distribution $G(z)$. With probability $G(z=x)$, renters paying $p_{r,1}h$ in the free market in the current period face additional rent payments of $x p_{r,1}h$ in the subsequent period, where $x$ represents an independent shock that can take negative as well as positive values. Consequently, the average per-unit rent in the free market is given by $\int_\mathcal{Z}(1+z)p_{r,1}\mathrm{d}G(z)$, denoted as $\overline{p}_{r,1}$.

In contrast, rent in the controlled market is certain and diminishes with tenancy duration. Similar to \citet{fmn-2023}, prolonged tenancy in the controlled market leads to lower real rent. Following the approach of \citet{fmn-2023}, controlled rent is determined by tenancy duration as follows:

\begin{equation}\label{eq:rent_in_controlled_market}
\frac{p_{r,0}}{(1+\kappa_{1})^{\iota}} h,
\end{equation} 

\noindent where $h$ denotes the controlled unit at the time a new tenant moves in. $p_{r,0}$ represents the initial rental price paid by new tenants, while $\kappa_{1}$ signifies the rate at which rent control diminishes real rents for each additional period tenants remain in the same unit. $\iota$ indicates the number of years they have resided in that unit.

Unlike landlords of free market units who fully offset the depreciation of their properties each period, landlords of controlled units neglect their units because of rent control. Consequently, controlled units deteriorate over time, resulting in a reduction of housing services available to tenants \citep{gl-1990,la-2018}. The amount of housing services tenants obtain in the controlled market is thus given by:

\begin{equation}
    h(o,f,\iota)=\frac{1}{(1+(1-o)(1-f)\kappa_{2})^{\iota}} h.
\end{equation}
 
\noindent $\kappa_{2}$ represents the rate at which controlled units deteriorate for each additional period a household remains renting the same unit.\footnote{We do not keep track of units that have depreciated under rent control (or any other units). When tenants in rent-controlled units vacate, these units can be subdivided or restructured in any manner for new renters.}

Figure \ref{fig:diagram_rent_path} provides a qualitative illustration of rents paid by renters as a function of the market in which they rent and their tenancy duration.\footnote{To maintain clarity in our illustration, we abstract from the effects of rental unit deterioration.} The downward sloping blue line depicts rents paid by renters residing in controlled units. Upon entering a new rental unit, their rent is $p_{r,0}$, but as their tenancy duration increases, their rent decreases in accordance with equation \eqref{eq:rent_in_controlled_market}.

Conversely, the horizontal red line represents the average rents paid by renters dwelling in the free market, which we denote as $\overline{p}_{r,1}$. This average rent remains constant irrespective of their tenancy duration, although their actual rent may fluctuate around this average value.

Figure \ref{fig:diagram_rent_path} indicates that in equilibrium, the average rent paid in the free market, $\overline{p}_{r,1}$, must be lower than the initial rent paid by renters in the controlled market, $p_{r,0}$. However, the average rent in the controlled market, $\overline{p}_{r,0}$, is lower than $\overline{p}_{r,1}$. What accounts for this equilibrium outcome?

We explain this equilibrium outcome as follows. If the initial rents paid by new tenants in the controlled market were lower than the average rents paid by renters in the free market, $\overline{p}_{r,1}$, all renters would opt to rent in the controlled market. To ensure market equilibrium, renters in the controlled market must initially pay a premium. Renters who anticipate remaining in the same unit for only a short duration and have sufficient risk tolerance for rent fluctuations would choose to rent in the free market to avoid paying this premium. Conversely, renters expecting to stay in the same unit for an extended period, and whose risk tolerance for rent volatility is low, would be willing to pay the premium to benefit from protection against rent fluctuations and lower long-term rents.

\subsection{Mortgages}\label{subsection:morgtages}

In addition to rent control, the availability of mortgage contracts significantly influences households' decisions to buy or rent and is another quantitatively important economic factor. Following \citet{kkk-2021}, we model mortgages, denoted as $m$, as long-term debt instruments with associated origination costs, prepayment penalties, and borrowing restrictions. We set the mortgage interest rate to $r_{m}=r+\Delta_m$, where $r$ represents the return on financial assets, and $\Delta_m$ signifies the mortgage debt interest rate premium \citep{sommeretal-2013,sommersullivan-2017,kaplanetal-2019}. Households cannot deduct mortgage interest payments from their taxable income because mortgage interest payments are not tax-deductible in Canada.

To obtain a new mortgage, households must incur a fixed cost denoted as $\omega_{1}$. At the time of a new mortgage origination only, the loan-to-value ratio (LTV) must not exceed a fraction $\lambda_{m,1}(ph)<1$, which depends on the value of the house at the time of purchase, denoted as $ph$:

\begin{equation}\label{eq:LTV}
\frac{m}{p h} \leq \lambda_{m,1}(ph).
\end{equation}

\noindent A constraint, which we refer to as the payment-to-income ratio (PTI), also imposes a ceiling on mortgage debt at origination:

\begin{equation}\label{eq:PTI}
\frac{\tau_{p}ph+(r_{m}+\upsilon_{1})m}{y_{\ell,j}(\chi,e)+y_{k}(a)} \leq \lambda_{m,2}.
\end{equation}

\noindent $\upsilon_{1}$ represents the fraction of debt paid every year on average over the life of a mortgage, while $\lambda_{m,2}$ denotes the maximum PTI ratio permitted. Constraint \eqref{eq:PTI} stipulates that property taxes, interest payments, and the average fraction of the principal paid cannot exceed a fraction $\lambda_{m,2}$ of their pre-tax labor and capital incomes.

Each period, a mortgage holder must allocate at least a fraction $\upsilon_{1}$ of their total mortgage debt towards repayment and is permitted to pay up to a fraction $\upsilon_{2}>\upsilon_{1}$ of the debt without incurring a penalty. However, any payment exceeding $\upsilon_{2}m$ is subject to a prepayment penalty, which we outline as follows:

\begin{equation}\label{eq:prepayment_penalty}
\tau_{pp}(m,m')=\omega_{2}\mbox{max}\left\{(1-\upsilon_{2})m-m',0\right\}.
\end{equation}

\noindent If a mortgage holder pays less than $\upsilon_{1}m$, we consider this akin to initiating a new mortgage. Consequently, they must incur the origination cost, $\omega_{1}$, and adhere to constraints \eqref{eq:LTV}--\eqref{eq:PTI} once again.

\subsection{Dynamic program}\label{subsection:household_optimization}
Given its state variables, a household must make decisions regarding various aspects such as saving ($a'$), homeownership or renting ($o' \in \left\{0,1\right\}$), cohabitation with their parents ($v' \in \left\{0,1\right\}$), or homelessness ($b' \in \left\{0,1\right\}$). Both homeowners and renters determine their housing consumption ($h'$). Renters additionally decide whether to rent in the controlled or free market ($f' \in \left\{0,1\right\}$, where we set $f'=1$ if $o'=1$, $v'=1$, or $b'=1$). Renters renting in the controlled market also choose whether to extend their tenancy by one period ($\iota'$). Homeowners decide on their mortgage debt level ($m'$) and whether to finance a new mortgage ($n$). Households reaching age $J$ also decide on the amount of bequests to leave behind ($q'$). We formulate the dynamic program encapsulating these decisions as follows:

\begin{align}
  \label{eq:vf}
  & V_{j}(s)=\max_{\{c,a',v',b',o',f',\iota',h',n,,m',q'\}}\ \biggl\{u_{j}\left(c,h'\right(o',f',\iota'\left)\right)+\indicator{j=J}\omega_{B}(q')+ \\ & \quad
  \indicator{j<J}\beta \int_{\mathcal{E}\times \mathcal{Z}}V_{j+1}(s')dF(e,e'|b')dG(z)\biggr\}. \notag 
\end{align}
Subject to
\begin{flalign}
  & c(1+\tau_{c})+a'+or_{m}m+(1-b')(1-v')(1-o')\left[(1-f')\frac{p_{r,0}}{(1+\kappa_{1})^{\iota'-1}}h'(0)+f'p_{r,1}h'\right]+\label{eq:bc} \\ & \quad \indicator{h>h_{p}}(1-o)fzp_{r,1}h+o\left(\tau_{p}+\delta_{h}+\indicator{o'=0 \vee h'\neq h}\tau_{s}\right)ph + o'(1+\tau_{b}\indicator{o=0 \vee h'\neq h})ph'+ \notag \\ & \quad\omega_{1}n+(1-n)\tau_{pp}(m,m') + \indicator{o'=1 \land j=J}ph(\tau_{p}-\delta_{h}-\tau_{s})= \notag \\ & \quad y_{\ell,j}(\chi,e)+y_{k}(a)-y^{\tau}_{j}(\chi,e,a)+a+o(ph-m)+o'm'+\vartheta(\chi,e) \notag \\
  & a' \geq 0 \label{eq:financial_wealth} \\
& h' \begin{cases}
\geq (1-b')(1-v') o' \underline{h}_{o}+(1-b')(1-v')(1-o')\underline{h}_{r} \hspace{0.5cm} \text{if $b' \neq 1$ and $v' \neq 1$.}\\
=b'\underline{h}_{h}+v'\underline{h}_{p} \hspace{0.5cm} \text{if $b' = 1$ or $v' =1$.}
\end{cases}\label{eq:min_housing} \\
  & v'=0 \hspace{0.5cm} \mbox{if } j\geq J_{p} \label{eq:can_live_with_parents}\\
  & b'=0 \hspace{0.5cm} \mbox{if } j<J_{p} \label{eq:can_be_homeless}\\
  & h'(o',f',\iota')=\frac{1}{(1+(1-o')(1-f')\kappa_{2})^{\iota'-1}}h' \label{eq:housing servies_controlled}\\
  & n \in \left\{0,o'\indicator{j<J_{R}}\right\} \label{eq:origination} \\
  & nm' \leq \lambda_{m,1}(ph') p h' \label{eq:LTV_1} \\
  & n (\tau_{p}ph'+(r_{m}+\upsilon_{1})m')\leq \lambda_{m,2}(y_{\ell,j}(\chi,e)+y_{k}(a))\label{eq:PTI_1} \\
  & (1-n)m'\leq (1-\upsilon_{1})m \label{eq:settle_debt} \\
  & q' = a' + o' \left[p h'-m'\right] \label{eq:networth}
\end{flalign} 

\noindent Household utility comprises three elements: a flow utility derived from housing and non-housing consumption, a warm-glow component stemming from bequests left behind by deceased households, and a discounted future value.\footnote{We formulate the dynamic program in a manner that accommodates labor market productivity shocks for retirees and working-age households, ensuring brevity in notation. This approach enables us to represent both retired and working-age households' using a single equation.}

Rentals units occupied by households in the controlled market deteriorate over time due to landlords' neglect to maintain them, resulting in a decline in the consumption of housing services.

Equation \eqref{eq:bc} delineates the household's budget constraint, encompassing various expenditures and available resources. On the left-hand side, expenditures include (in order of appearance): consumption and consumption taxes, savings, interest payments for mortgage holders, costs of renting in both controlled and free markets, rent shocks, property taxes, depreciation costs, home selling costs, home purchasing costs, origination costs for new mortgages, prepayment penalty costs for mortgage holders, and property taxes, depreciation, and selling costs for households who decide to die as homeowners. Conversely, the right-hand side comprises resources available for expenditure, which we list in the order of their appearance: after-tax labor and capital incomes, financial assets, housing equity for owners, mortgage funds raised, and bequests contingent upon the household's labor productivity. This budget constraint framework provides a comprehensive overview of the financial dynamics underlying household decision-making processes.

Constraint \eqref{eq:financial_wealth} stipulates that households lack access to unsecured debt. Constraint \eqref{eq:min_housing} mandates that all owners must consume at least $\underline{h}_{o}$ units of housing, while all renters must consume at least $\underline{h}_{r}$ units. It also says that homeless households must consume $\underline{h}_{h}$ units of housing, and households cohabiting with their parents must consume $\underline{h}_{p}$ units of housing. Constraint \eqref{eq:can_live_with_parents} specifies that households whose age is greater than or equal to $J_{p}$ cannot reside with their parents. Constraint \eqref{eq:can_be_homeless} states that households whose age is strictly less than $J_{p}$ can live with their parents and therefore cannot be homeless.\footnote{As highlighted in the introduction, research conducted by \citet{spatial_distance-2020} indicates that around 75\% of young adults in the United States, aged between 26 and 34, live within a 30-mile radius (roughly 48 kilometers) of their parental residence. This statistic implies that a significant portion of young adults in Toronto likely have the opportunity to reside with their parents if they choose to do so.}

Constraint \eqref{eq:housing servies_controlled} says that renters in the controlled market experience reduced housing services consumption as they extend their stay in the same unit. Constraint \eqref{eq:origination} asserts that a household can only originate a new mortgage if it is of working age. Constraints \eqref{eq:LTV_1} and \eqref{eq:PTI_1} represent the LTV and PTI constraints at mortgage origination, respectively. Constraint \eqref{eq:settle_debt} mandates that mortgage holders must repay at least a fraction $\upsilon_{1}$ of their mortgage principal. Lastly, constraint \eqref{eq:networth} equates the bequests left by deceased households to the sum of their financial and net housing wealth.

We denote the policy functions for consumption, savings, living with parents, becoming homeless, housing tenure, rental market choice, years living in the same unit, housing size, mortgage debt, mortgage origination, and bequests by $c_{j}(s)$, $a'_{j}(s)$, $b'_{j}(s)$, $v'_{j}(s)$, $o'_{j}(s)$, $f'_{j}(s)$, $\iota'_{j}(s)$, $h'_{j}(s)$, $m'_{j}(s)$, $n_{j}(s)$, and $q'_{j}(s)$, respectively.

\subsection{Rental company}\label{sec:rental_company}

We adopt the rental housing supply theory outlined in \citet{fallis-1985} and incorporate it within the rental housing supply framework developed by \citet{garrigaeetal-2009}. In our formulation, we designate $i=0$ to represent the controlled market and $i=1$ to represent the free market. Within each market $i$, for every rental housing unit, denoted by $x_i = 1$, the rental company obtains an average rent, $\overline{p}_{r,i}$, while also bearing property taxes, $\tau_{p}$, and depreciation costs, $\delta_{h}$.

In order to provide a total of $X'=\sum^{1}_{i=0} X'_{i}$ rental housing, the rental company also faces management expenses amounting to $\frac{\theta_{1}}{\theta_{2}}X'^{\theta_{2}}$. $\theta_{1}$ signifies the average cost associated with managing the rental housing stock $X'$, while $\theta_{2}\geq 1$ delineates the price elasticity of the rental housing supply.

The rental company is mandated by the government to provide a minimum quantity of controlled rental housing, which we denote as $\mbox{min}\{\underline{X}_{0},X'\}$, where $\underline{X}_{0} \geq 0$. If the total quantity of rental housing supplied, $X'$, is equal to or less than $\underline{X}_{0}$, the market falls entirely under control. In such instances, the rental company's marginal revenue is the average rent prevailing in the controlled market, denoted as $\overline{p}_{r,0}$. Conversely, if $X'$ exceeds $\underline{X}_{0}$, the rental company's marginal revenue becomes the average free-market rent, denoted as $\overline{p}_{r,1}$. We use the term ``marginal rent" interchangeably with the rental company's marginal revenue. In scenarios in which $X'$ surpasses $\underline{X}_{0}$, $\underline{X}_{0}$ signifies the quantity of controlled units supplied, while $X' - \underline{X}_{0}$ represents the quantity of rental housing supplied in the free market.

The rental company purchases houses at the prevailing house price, $p$. In each period, it determines the quantity of rental housing to provide in the controlled market, represented by $X'_{0}$, and the quantity to supply in the free market, denoted as $X'_{1}$. The rental company is confronted with the following optimization problem:

\begin{align}\label{eq:rental_company_problem}
V(X)=\underset{X'_{0},X'_{1}} {\text{$\mbox{max} $}}\left\{\sum_{i=0}^{1} \overline{p}_{r,i}X'_{i}-\frac{\theta_{1}}{\theta_{2}}X'^{\theta_{2}}-p X'+p(1-\delta_{h}-\tau_{p})X+\frac{1}{1+r}V(X')\right\}.
\end{align}
Subject to
\begin{align}
& X'_{0}=\mbox{min}\{\underline{X}_{0},X'\} \label{eq:supply_restriction} \\
& X'_{1} \geq 0 \label{eq:supply_positive} \\
& X'=\sum^{1}_{i=0} X'_{i}. \label{eq:total_supply} 
\end{align}

\noindent Equation \eqref{eq:supply_restriction} stipulates that the quantity of rental housing supplied in the controlled market is the lesser of $\underline{X}_{0}$ and $X'$. Equation \eqref{eq:supply_positive} ensures that the quantity supplied in the free market is non-negative. Equation \eqref{eq:total_supply} denotes that the total quantity supplied is the sum of the quantities supplied in both markets. When $X'>\underline{X}_{0}$, the supply function is expressed as follows:

\begin{align}\label{eq:rental_supply}
X'^{\star}=\left[\frac{1}{\theta_{1}}\left(\overline{p}_{r,1}-p\left(\frac{r+\delta_{h}+\tau_{p}}{1+r}\right)\right)\right]^{\frac{1}{\theta_{2}-1}}.
\end{align}

\noindent In equilibrium, when $X'^{\star} \geq \underline{X}_{0}$, it follows that $X'^{\star}_{0}=\underline{X}_{0}$ and $X'^{\star}_{1}=X'^{\star}-\underline{X}_{0}$. However, if $X'^{\star} < \underline{X}_{0}$ or if rent control laws are applicable to all rental units, then in equation \eqref{eq:rental_supply}, $\overline{p}_{r,0}$ replaces $\overline{p}_{r,1}$, and $X'^{\star}_{0}=X'^{\star}$ and $X'^{\star}_{1}=0$.

\noindent The partial derivatives of $X'^{\star}$ with respect to $p$ and $\overline{p}_{r,1}$ reveal that the optimal quantity supplied:

\begin{itemize}
    \item Decreases with the housing price and property tax rate, and increases with $\overline{p}_{r,1}$ if $\theta_{2} > 1$, as is the case in our calibration.
    \item Remains independent of $\overline{p}_{r,1}$ if $\theta_{2} = 1$; in this scenario, rental supply is infinitely elastic.
\end{itemize}

One might contend that over time, the proportion of the controlled rental market naturally diminishes due to the depreciation of existing rent-controlled units, which may be supplanted by new constructions not bound by rent control regulations. However, to sustain a consistent share of the rent-controlled market in a stationary equilibrium, governments can raise the threshold year of new construction beyond which units are exempt from rent control. This approach mirrors the actions taken by the Ontario government over the years, and is similar to measures implemented in Quebec, where newly constructed units are granted a five-year exemption from rent control regulations.

\subsection{Housing construction company}\label{sec:construction_company}

A competitive construction firm engages in the construction of new housing \citep{rotberg-2022}. It earns revenues, $pI$, from its production of new housing, denoted by $I$. However, the firm also bears costs, which are encapsulated by the following technology:

\begin{equation*}
TC(I)=\frac{\epsilon_{1}}{\epsilon_{2}}I^{\epsilon_{2}}\label{eq:total_cost_structure},
\end{equation*}

\noindent where $\epsilon_{1}$ dictates the average cost of housing investment, while $\epsilon_{2}$ governs the price elasticity of housing investment. The construction company makes decisions regarding investment to maximize its profits:

\begin{equation*}
\label{eq:construction_company_profits}
\underset{I} {\text{$\mbox{max}$}}{\left\{p I-TC(I)\right\}}.
\end{equation*}

\noindent The solution yields the following relationship between housing investment and the price of housing:

\begin{equation}
\label{eq:construction_company_first_order_condition}
I^{\star}=\epsilon_{1} p^{\epsilon_{2}-1}.
\end{equation}

\subsection{Aggregation}\label{subsection:aggregation}

As we outlined in the general description of the model, we incorporate both federal and local governmental entities. The federal government levies income and consumption taxes while also providing pensions to households. On the other hand, local governments collect property taxes and supply public goods to residents and shelters services to households experiencing homelessness. Given that rent-control policies are typically set at the local level, our focus is on ensuring budgetary balance at the local government level:

\begin{align}\label{eq:gbc}
  G + \sum_{j=1}^J \int_\mathcal{S} \theta_{b} b'_{j}(s)\mathrm{d}\Psi(s)  = \tau_{p} p H.
\end{align}

\noindent We require that government expenditures on public goods and support for homeless households be matched by revenues from property taxes. When calibrating the model to align with the current tax code in Toronto, we determine the level of consumption of public goods, which we denote as $G$, to achieve a balanced budget. Throughout our experiments, we maintain the value of $G$ constant and adjust $\tau_{p}$, the property tax rate, to ensure budget balance.%\footnote{Federal taxes fund pensions and public goods for households across the country. Since labor supply is inelastic, labor income tax revenues remain unchanged across equilibria. Although tiny fluctuations in aggregate capital income and consumption affect federal tax revenues coming from the city's economy, the impact on city residents through lump-sum transfers would be essentially null, as the transfers would be distributed to households nationwide. In our simulations, such transfers would be essentially zero. Therefore, we abstract from balancing the federal budget constraint.}

$S=A\times O \times F \times I \times Z \times H\times M\times X \times E$ denotes a typical subset of the state space $\mathcal{S}$; and $A$, $O$, $F$, $I$, $Z$, $H$, $M$, $X$ $E$ denote typical subsets of $\mathcal{A}$, $\mathcal{O}$, $\mathcal{F}$, $\mathcal{I}$, $\mathcal{Z}$, $\mathcal{H}$, $\mathcal{M}$, $\mathcal{X}$, and $\mathcal{E}$, respectively.

Aggregate market housing demand must equal aggregate housing supply:

\begin{equation}
\label{eq:housing_clearing_condition}
  \sum_{j=1}^J \int_\mathcal{S} (1-b'_{j}(s))(1-v'_{j}(s)) \left(\frac{h'_{j}(s)}{(1+(1-o'_{j}(s))(1-f'_{j}(s))\kappa_{2})^{\iota'_{j}(s)}}\right) \mathrm{d}\Psi(s) = H',
\end{equation}

\noindent and the housing stock law of motion is: 

\begin{equation}\label{eq:housing_stock_lom}
H'=I+(1-\delta_h)H.
\end{equation}

\noindent Equation \eqref{eq:housing_stock_lom} represents a standard law of motion for housing, stating that the next period's housing stock equals housing investment plus the depreciated stock of housing. In each rental market, the aggregate rental housing demand must match the aggregate rental housing supply:

\begin{equation}\label{eq:rental_housing_clearing_condition}
  \sum_{j=1}^J \int_\mathcal{S} (1-b'_{j}(s))(1-v'_{j}(s)) (1-o'_{j}(s))\indicator{f'_{j}(s)=i}\left(\frac{h'_{j}(s)}{(1+(1-f'_{j}(s))\kappa_{2})^{\iota'_{j}(s)}}\right) \mathrm{d}\Psi(s) = X'_{i} \quad i \in \{0,1\}.
\end{equation}

The distributions of households evolve according to the laws of motion,
\begin{equation}\label{eq:lom1}
\Psi_{1}(\mathrm{S})=\int_{\mathcal{X} \times \mathcal{E}} \indicator{(0,0,0,0,0,\underline{h}_{r},0,X,E)\in \mathrm{S}}\mathrm{d} \hat{F}(X)\tilde{F}(E),
\end{equation}


\begin{equation}\label{eq:lom2}
\Psi_{j+1}(S)=\int_{\mathcal{S}} \hat{F}(X)F(E,e|b'_{j}(s))G(z)T_{j}(s)\mathrm{d}\Psi_{j}(s),\ j=1,\ldots,J-1
\end{equation}

\noindent where,

\begin{align}
    T_{j}(s) = \indicator{[a'_{j}(s)\in A] \land [o'_{j}(s)\in O] \land [f'_{j}(s)\in F] \land [\iota'_{j}(s)\in I] \land [h'_{j}(s) \in H] \land [m'_{j}(s) \in M]}.\notag
\end{align}

Finally, total bequests left by decedents must equal the total wealth endowments:
\begin{equation}\label{lom:agg_bequests}
      \int_{\mathcal{S}} q'_{J}(s) \mathrm{d}\Psi_{J}(s) =  \sum_{j=1}^J \int_{\mathcal{S}} \vartheta (\chi,e)\ \mathrm{d}\Psi_j(s).
\end{equation}

\noindent Similar to the methodology employed by \citet{kkk-2021}, we select the parameter $\vartheta$ to ensure this condition holds across all our experiments. Consequently, variations in $\vartheta$ between equilibria are contingent upon proportional changes in aggregate bequests.

\subsection{Equilibrium}
An equilibrium in the model is a sequence of aggregate prices and quantities, $\left\{H',I,\{X'_{i},\overline{p}_{r,i}\}^{1}_{i=0},p_{r,0},p,\vartheta \right\}$, a sequence of value and policy functions, $\biggl\{V_{j}(\cdot),c_{j}(\cdot),a'_{j}(\cdot),v'_{j}(\cdot),b'_{j}(\cdot),o'_{j}(\cdot),f'_{j}(\cdot),\iota'_{j}(\cdot),h'_{j}(\cdot),m'_{j}(\cdot),n_{j}(\cdot),q'_{j}(\cdot)\biggl\}_{j=1}^J$, and a sequence of measures, $\left\{\Psi_{j}(\cdot)\right\}_{j=1}^J$, that:
\begin{enumerate}
\item Solve the household's dynamic problem (\ref{eq:vf}) subject to constraints (\ref{eq:bc})--(\ref{eq:networth});
\item Solve the rental company's problem (\ref{eq:rental_company_problem}) subject to (\ref{eq:supply_restriction})--(\ref{eq:total_supply});
\item Satisfy the construction company's first-order condition, (\ref{eq:construction_company_first_order_condition});
\item Satisfy the government's budget constraint, (\ref{eq:gbc});
\item Clear the housing market,
(\ref{eq:housing_clearing_condition}) subject to the housing stock law of motion (\ref{eq:housing_stock_lom});
\item Clear both rental housing markets,
(\ref{eq:rental_housing_clearing_condition});
\item Satisfy the laws of motion for the distribution of households, (\ref{eq:lom1})--(\ref{eq:lom2});
\item And satisfy the law of motion for bequests, (\ref{lom:agg_bequests}).
\end{enumerate}

\section{Calibration and model performance}\label{sec:calibration}

Our calibration involves two steps. Initially, we calibrate specific parameters using data from Toronto and estimates from the literature. Then, we calibrate the values of the remaining parameters to align the model's generated statistics with certain moments of Toronto data. Table \ref{tab:calibration_external} details the parameter values we set externally along with their sources. Similarly, Table \ref{tab:calibration_internal} outlines the parameter values we set internally alongside their targets. In Section \ref{appendix:data} of the online appendix, we elaborate on how we derived certain statistics using various datasets. Lastly, in Section \ref{subsection:non_targeted_statistics}, we demonstrate the model's performance across several dimensions we did not target.

\subsection{Assigned parameter values}

\subsubsection{Demographics, preferences, and income}
A model period spans 1 year. Households live from age 26 to age 85, and retire at age 66, which implies a retirement age of $J_{R}=41$ and a lifespan of $J=60$. Following \citet{rs-2023}, we set the maximum age at which households can live with their parents, $J_{p}$, to 9, thereby prohibiting households aged 35 and above from living with their parents. We set the households' relative risk aversion at $\sigma=2$, and establish the equivalent household size for each age group, $\{\xi_{j}\}_{j=1}^{J}$, using 2016 Census data on the average household size for each age cohort. We scale household sizes by their square roots, following standard practice in the literature. We determine the discount factor, $\beta$, and the utility weight on housing services, $\gamma$, in the second stage of our calibration.

We adopt the labor productivity process outlined in \citet{gkkoc}. We set the standard deviation of the fixed effect to $0.128$, and the stochastic labor productivity shocks follow an AR(1) process in logarithmic terms,

\begin{equation*}
  \log e' = \rho_{e} \log e + \epsilon_{e},\quad \epsilon_{e}\overset{\mathrm{i.i.d.}} \sim N(0,\sigma_{e}).
\end{equation*}

\noindent We set the persistence parameter, $\rho_{e}$, to 0.9, following \citet{gkkoc}. In the second stage, we determine the standard deviation of the stochastic component of labor productivity, denoted as $\sigma_{e}$. Newborns obtain their stochastic labor productivities from the unconditional probability distribution associated with the ergodic distribution, which is endogenously generated by our productivity process. Consistent with \citet{az-2024}, for households who opt for or are compelled to be homeless, we mechanically increase the probability of remaining in their stochastic labor productivity state, $e$, by 20\% compared to what is implied by our aforementioned stochastic labor productivity process.

Based on the 2016 Census data, we establish the deterministic lifecycle component of labor market productivity as follows:

\begin{equation*}
\zeta_j = 0.923+0.09\times j-1.61\times 10^{-3}\times j^{2}, \mbox{ for } j\leq J_{R}
\end{equation*}

\noindent We derive these estimates by regressing the average household income for each age, ranging from 26 to 64, on age and age squared. We determine pension benefits, denoted as ${y_{R}(\chi,e)}$, based on the replacement rate of gross individual earnings, as outlined in \citet{OECD-pension}.

Surviving households' endowments are determined proportionally to their labor income, expressed as $\vartheta(\chi,e):=\vartheta \chi e$. We select the parameter $\vartheta$, governing surviving households' wealth endowments, to fulfill equation \eqref{lom:agg_bequests}.

\subsubsection{Taxes and interest rates}

We set the taxable portion of capital income, $\nu$, to be 50\%, mirroring the provision in Canada's tax code. Utilizing household-level data from the 2016 Census for Toronto, we estimate $\psi=0.9275$ in equation \eqref{eq:income_tax}. We set the consumption tax rate, denoted as $\tau_{c}$, to 13\% to align with the Harmonized Sales Tax (HST) rate in Ontario. Furthermore, we establish the interest rate, $r$, at 2.77\%, as estimated for Canada by \citet{bruneau-etal-2018}.

\subsubsection{Housing}

We set our rent shocks $\left\{z\right\}_{\mathbb{Z}}$ to $\left[-0.1227,-0.0607,0,0.0624,0.1284,0.2589\right]$, along with corresponding probabilities $\{G(z)\}_{\mathbb{Z}}$ given by $\left[0.16,0.12,0.44,0.08,0.16,0.04\right]$. We derive these values from 1992 to 2016 data sourced from the Rental Market Survey (RMS), reflecting average rents in both controlled and free markets. Our analysis indicates that in Toronto, there is a non-negligible higher risk of substantial rent increases relative to the chance of significant rent decreases.

To determine the rate at which real rents decrease with tenancy duration in the controlled market, denoted as $\kappa_{1}$, we refer to 2016 Census data, setting it at 1.365\%. Similarly, we set $\kappa_{2}$, representing the annual deterioration rate of controlled units, to 0.55\%, using estimates from \citet{murray_et_al-1991} for Los Angeles, which at the time followed a policy closely aligned with tenancy control as we model here. This figure aligns well with the economic notion that landlords, compelled to provide real rent reductions each additional period a tenant remains in their unit, would opt to maximize their profits by allowing the property to deteriorate. Given that we set the housing depreciation rate, denoted as $\delta_{h}$, to 0.6\% based on estimates by \citet{rosenthal-2014}, it seems sensible to employ 0.55\% for $\kappa_{2}$.

We establish the per-household cost borne by the city to support the homeless, denoted as $\theta_{b}$, at 82.3\% of the average household income. We determine this parameter by dividing the City of Toronto's total spending on homelessness supports by the number of homeless individuals in the city in 2021.\footnote{See \url{https://www.homelesshub.ca/community-profile/toronto} and \url{https://www.toronto.ca/legdocs/mmis/2021/ex/bgrd/backgroundfile-163804.pdf}.}.

We set the property tax rate, denoted as $\tau_p$, to 0.44\%, drawing upon information provided by Altus Group and 1992-2016 Property Sale and Assessment Data (PSAD). Following \citet{gm-2003}, we determine buying costs, $\tau_{b}$, as 2\% of the purchase price, and selling costs, $\tau_{s}$, as 6\% of the sale price.

To achieve a price elasticity of housing supply equal to 0.89, as estimated by \citet{piaxao-2021} for Toronto, we select $\epsilon_{2}$ to be 1.89. Additionally, the construction company's cost function level, denoted as $\epsilon_1$, is a free parameter that we derive from the market clearing condition \eqref{eq:housing_clearing_condition}.

Lastly, we set the weighted average rent in both rental markets, denoted as $\overline{p}_{r}$, to 5.24\% of the equilibrium house price. We base our computation on 2016 data we obtained from the Canadian Real Estate Association (CREA).

\subsubsection{Mortgages}

In accordance with banking regulations in Canada, we set the loan-to-value (LTV) constraint, denoted as $\lambda_{m,1}(ph)$, as follows:

\begin{equation*}
\lambda_{m,1}(ph) = \begin{cases}
0.8 &\text{if $ph\geq \$1M$}\\
0.9 &\text{if $\$500K\leq ph< \$1M$}\\
0.95 &\text{if $ph < \$500K$}.
\end{cases}
\end{equation*}

\noindent Mortgage holders with a model-equivalent house value above \$1M can borrow only up to 80\% of its value; mortgage holders with a model-equivalent house value between \$500,000--\$1M can borrow up to 90\% of the house value; and mortgage holders with a house value below \$500,000 can get a loan of up to 95\% of the house value. We set $\lambda_{m,2}=0.39$ so that the maximum allowable PTI is 39\% as is standard in the mortgage industry in Canada.\footnote{See \url{https://www.cmhc-schl.gc.ca/consumers/home-buying/calculators/debt-service-calculator}.}

For mortgage contracts, we begin by setting the fixed mortgage origination cost, denoted as $\omega_{1}$, to 0.68\% of the average household income in Toronto. We determine this based on data we source from the 2016 Census and the Canada Mortgage and Housing Corporation's (CMHC) website, which includes expenses such as appraisal fees (approximately \$250) and lawyer fees (approximately \$500).\footnote{See \url{https://www.cmhc-schl.gc.ca/en/professionals/industry-innovation-and-leadership/industry-expertise/resources-for-mortgage-professionals/how-much-will-my-home-really-cost}}

We set the prepayment penalty, denoted as $\omega_{2}$, to 3 months of interest, a standard practice in Canadian mortgage contracts. Following the methodology in \citet{kkk-2021}, we set $\upsilon_{1}$, which represents the average fraction of mortgages paid each year, to 1.7\% to mirror a typical 30-year amortization period prevalent in the Canadian mortgage industry. We establish the threshold for the prepayment penalty, $\upsilon_{2}$, at $1.15 \times \upsilon_{1}$ plus 15\% of the total debt and interest payments, amounting to 18.8\%. This threshold is aligned with standard mortgage contracts in Canada, allowing prepayments of 15\% of the principal and regular payments without penalty.

Lastly, we set the interest rate premium on mortgage debt, denoted as $\Delta_{m}$, to 0.19\%. This figure represents the average spread between the 5-year variable rate on mortgages and the return on a 10-year government bond, spanning from 2002 to 2022.\footnote{The mortgage data we use is not publicly available. The 10-year Government of Canada bond yield can be found at \url{https://www.bankofcanada.ca/rates/interest-rates/canadian-bonds/}}.

\subsection{Internally calibrated parameters}

After assigning the parameter values mentioned earlier, we proceed to calibrate another twelve parameters, as we outline in Table \ref{tab:calibration_internal}. These variables encompass the discount factor, $\beta$; utility weight on housing services consumption, $\gamma$; strength of bequest motive, $\mu$; average income taxation level, $\tau_{\ell}$; variability in labor income's stochastic component, $\sigma_{e}$; convexity of rental housing management, $\theta_{2}$; marginal managerial cost-to-rent ratio determined by $\theta_{1}$; initial rent paid by a tenant entering a unit under tenancy control, $p_{r,0}$; minimum rental unit size, $\underline{h}_{r}$; minimum housing unit size for purchase, $\underline{h}_{o}$; housing services for households living with their parents, $\underline{h}_{p}$; and housing services for homeless households, $\underline{h}_{h}$.\footnote{The construction cost scalar, $\epsilon_{1}$, is a free parameter we choose to equate housing supply to housing demand in the benchmark equilibrium. In counterfactual simulations, we maintain $\epsilon_{1}$ at its benchmark level.}

We collectively calibrate all these parameters to align with twelve key metrics in Toronto data:

\begin{enumerate}
    \item The ratio of aggregate mortgage debt to aggregate house value among homeowners, which stands at 21.5\% according to the 2019 Survey of Financial Security (SFS).
    \item The proportion of total renters' income dedicated to rent, which we estimate at 22.5\% using 2016 Census data.
    \item The ratio of aggregate bequests to aggregate net wealth, which we calculate at 2.2\%, utilizing estimates we derive from the 2019 SFS.
    \item The average labor income tax rate, recorded at 19.6\%, which we source from the 2016 Census.
    \item The Gini coefficient of labor income, which we also obtain from the 2016 Census.
    \item The property tax burden on renters in cities akin to Toronto, which we estimate at approximately 12\%. We ascertain this figure by averaging estimates for similar cities from multiple sources, including \citet{cy-1994}, \citet{susin-2002}, \citet{fack-2006}, and \citet{sy-2020}.
    \item The ratio of average free-market rent to house price, which we calculate at 5.82\%, based on our analyses utilizing data from the 2016 Census and 2016 CREA.
    \item The percentage of renters in the controlled rental market, which we compute at 71\% using the 2016 Census. 
    \item The percentage of renters spending over 50\% of their before-tax income on rent, standing at 15.4\% according to the 2016 Census.
    \item The homeownership rate, reported at 66.5\% in the 2016 Census.
    \item The percentage of 26-34 year old households living with their parents, which we estimate at 22.3\% based on the 2016 Census.
    \item The percentage of households experiencing homelessness, which is documented at 0.26\% according to the ``Homeless Hub".\footnote{See \url{https://www.homelesshub.ca/community-profile/toronto}.}
\end{enumerate}

While we determine all parameters jointly, each parameter tends to exert a more significant influence on specific outcomes. Here is how:
\begin{enumerate}
    \item Discount factor, $\beta$: Primarily shapes homeowners' debt levels; a lower $\beta$ leads to higher homeowner debt levels as it influences their optimal borrowing behavior.
    \item Utility weight on housing, $\gamma$: Influences the proportion of renters' income spent on housing; higher $\gamma$ values indicate a greater desire for housing services among households.
    \item Bequest motive strength, $\mu$: Affects the ratio of bequests to net wealth; higher $\mu$ values result in relatively more bequests compared to overall wealth.
    \item Average income taxation level, $\tau_{\ell}$: Determines the average income tax rate; lower $\tau_{\ell}$ values correspond to higher tax burdens on households' incomes.
    \item Variability in labor income's stochastic component, $\sigma_{e}$: Shapes labor income inequality; higher $\sigma_{e}$ values lead to a greater Gini coefficient for the labor income distribution.
    \item Marginal cost of managing rental housing, $\theta_{2}$: Dictates the elasticity of rental housing supply; higher $\theta_{2}$ values result in less elastic rental supply and lower property tax burdens for renters.
    \item Marginal managerial cost-to-rent ratio, $\theta_{1}$: Influences the average free-market rent-to-price ratio; higher ratios are associated with higher $\theta_{1}$ values.
    \item Initial rent for controlled units, $p_{r,0}$: Controls the proportion of renters in the free market; higher initial rents relative to the free market lead to a larger share of renters opting for the free market.
    \item Minimum rental unit size, $\underline{h}_{r}$: Governs the proportion of renters spending over 50\% of their before-tax income on housing; larger minimum sizes result in a higher percentage of renters facing high housing cost burdens.
    \item Minimum housing unit size for purchase, $\underline{h}_{o}$: Influences the homeownership rate; larger minimum sizes lead to fewer households opting for homeownership.
    \item Housing services for households living with their parents, $\underline{h}_{p}$: Impacts the percentage of 26-34 year old households living with their parents; higher services lead to a higher proportion of households residing with their parents.
    \item Housing services for homeless households, $\underline{h}_{h}$: Affects the homelessness rate; greater services result in a higher rate of homelessness.
\end{enumerate}

\subsection{Model performance}\label{subsection:non_targeted_statistics}

We assess the performance of our model by examining aggregate, lifecycle, and tenure distribution statistics. Table \ref{tab:aggregate_non_targreted_statistics} presents the model's performance on aggregate statistics we did not target in our calibration.

Our model indicates a slightly higher ratio of average owned-house value to average household income compared to the data, alongside a somewhat lower percentage of owners with positive mortgage debt. Notably, our available data does not allow direct estimation of the ratio of average rental unit value to average owner-occupied unit value in Toronto. Hence, we compare our model's estimate with that of \citet{chatterjee-eyigungor-2015} for the United States, revealing a slight discrepancy.

Although directly comparable estimates are absent in our dataset, we include one additional statistic for reference. Our model projects average free-market rents exceeding average controlled rents by 3\%. Moreover, renters in the free market earn around 2.7\% more than their counterparts in the controlled market, a figure lower than observed in the 2016 Census.

The wealth distribution depicted in our model exhibits less wealth concentration at the top of the distribution compared to real-world data. This divergence arises primarily from the absence of rate-of-return heterogeneity, a factor that \citet{fagereng} highlight as pivotal for driving wealth concentration at the upper echelons of the distribution. Our model also lacks features such as heterogeneous rate of time preferences \citep{ks-1998} or income shocks akin to those observed in studies such as \citet{mj-shocks}, which could further elucidate our inability to replicate high levels of wealth concentration.

Figure \ref{fig:life_cycle} illustrates the lifecycle outcomes of the model in comparison to the data. In panel (a), we observe that the homeownership rate across the lifecycle in our model is lower than the data for both young and old age groups, while slightly higher for middle-aged households. Moving to panel (b), young households in our model exhibit higher LTV ratios compared to their counterparts in the data, whereas the opposite holds true for older households.

Figure \ref{fig:rent_to_income_cdf} displays the cumulative frequency distribution (CDF) of rent-to-income ratios. Our model closely mirrors the rent-to-income ratios observed in the data, albeit with a slight lower number of households with high rent-to-income ratios and a somewhat higher number of households with low rent-to-income ratios.

Table \ref{tab:tenure_by_income_quintile} presents households' tenure categorized by their income quintiles. The second column from the left corresponds to the lowest income quintile, while the first column from the right represents the highest income quintile. We disaggregate tenure into two categories: the percentage of renters in each quintile residing in controlled rental units and the percentage residing in free market units. Panel (a) displays these statistics based on the data, while panel (b) showcases the calibrated model's performance regarding the same metrics. Although the calibrated model does not perfectly replicate the data, it captures significant trends: notably, households in lower income quintiles are more inclined to rent in the controlled market compared to those in higher income quintiles.

\section{Experiments}\label{section:experiments}

\paragraph{Experiments' overview} Using our baseline calibrated model, we conduct four experiments.

\begin{itemize}
    \item In our first experiment, we compare the benchmark of our baseline model, which incorporates tenancy control alongside a free market, with the long-run equilibrium of our adjusted model featuring unit control instead of tenancy control.
    \item In the second experiment, we contrast the benchmark with the long-run equilibrium in which there is no rent control, which we call full decontrol.
    \item For the third experiment, we compare the baseline benchmark with the long-run equilibrium in which all rental housing units are subject to tenancy control, which we refer to as full tenancy control.
    \item In the last experiment, we search for the degree of tenancy control that maximizes newborns' welfare in the long run.
\end{itemize}

Across all experiments, we keep the tax code fixed and balance the government's budget by adjusting the property tax rate.

\paragraph{Welfare measure} Our primary welfare measure is the expected ex-ante utility of newborns \citep{ralws-1971,conesa-et-al,gkkoc}. This measure is akin to providing all newborns with a fixed proportional non-housing consumption transfer. Formally, the non-housing consumption-equivalent welfare effect resulting from a policy change is expressed as:

\begin{equation}
  CE_{t} = \left[\frac{\sum_{j=1}^{k}\int_{\mathcal{S}}V^{\dagger}_{j}(s)\mathrm{d}\Psi^{\dagger}_{j}(s)-\sum_{j=1}^{k}\int_{\mathcal{S}}\omega_{B}(q'^{\star}_{J}(s))\mathrm{d}\Psi^{\dagger}_{j}(s)}{\sum_{j=1}^{k}\int_{\mathcal{S}}V^{\star}_{j}(s)\mathrm{d}\Psi^{\star}_{j}(s)-\sum_{j=1}^{k}\int_{\mathcal{S}}\omega_{B}(q'^{\star}_{j}(s))\mathrm{d}\Psi^{\star}_{j}(s)}\right]^{\frac{1}{(1-\gamma)(1-\sigma)}}-1,
\end{equation}

\noindent where daggers denote counterfactual equilibrium objects, while stars denote objects in the benchmark. 

Additionally, we gauge the approval rate of the new policy change by quantifying the proportion of households experiencing an increase in their welfare under the new policy:

\begin{equation}
A = \sum_{j=1}^{k}\int_{\mathcal{S}} \indicator{V^{\dagger}_{j}(s)>V^{\star}_{j}(s)}\ \mathrm{d}\Psi^{\star}_{j}(s).
\end{equation}

\noindent For our primary welfare and approval measures, we set $k=1$, representing the welfare changes and approval rates of newborns. Additionally, we calculate average welfare changes and approval rates across all households by setting $k=J$. Since we compute approval rates based on the benchmark's distribution, we can further assess the approval among homeowners and renters in the benchmark.

\paragraph{Measures of rent} Our analysis encompasses several measures of rent that we outline below:

\begin{itemize}
    \item Average free-market rent ($\overline{p}_{r,1}$): The equilibrium average rent prevailing in the free market.
    \item Average controlled rent ($\overline{p}_{r,0}$): The equilibrium average rent observed in the controlled market.
    \item Average rent ($\overline{p}_{r}$): The weighted average of equilibrium rents across both the controlled and free rental markets.
    \item Rent of new tenants ($p_{r,0}$ in the controlled market and $\overline{p}_{r,1}$ in the free market): The equilibrium initial rent paid by new tenants upon entering a new rental unit. In the controlled market, new tenants face higher rents compared to long-standing tenants protected by rent control. In the free market, new tenants anticipate paying the same average rent as existing tenants in the free market.
\end{itemize}

\subsection{Tenancy control vs. unit control}

In our first experiment, we compare the benchmark with the long-run equilibrium featuring unit control. Before delving into the specifics of the model with unit control, we outline how we configure the unit control exercise to ensure comparability with the benchmark.

Firstly, we set the exogenously determined controlled rent of a unit in the unit control exercise to match the average controlled rent observed in our benchmark. Secondly, we align the proportion of the rental market under control in the unit control exercise with that of the benchmark. Lastly, to maintain consistency with the benchmark while ensuring a constant real rent that is independent of tenancy duration in the unit control exercise, we set the real rent decrease, $\kappa_{1}$ in the unit control exercise to equal the unit's deterioration rate, $\kappa_{2}$.

\subsubsection{Formal description of model with unit control}

We formalize the unit control model as follows: Each period, households ``wake-up" to encounter a random shock, which we denote as $\eta \in \{0,1\}$. The realization of $\eta$ yields a binary outcome: with probability $\pi$, $\eta$ equals 1; otherwise, it equals 0. Importantly, $\pi$ remains independent across time and previous realizations of $\eta$.

Households drawing $\eta = 1$ have the opportunity to rent in the controlled rental market in that period. Once in the controlled market, they retain the choice to stay in their current unit irrespective of subsequent realizations of $\eta$. However, upon vacating their unit, re-entry into the controlled market necessitates another realization of $\eta$ equal to 1. New tenants entering the controlled market initially pay a rent equivalent to the benchmark average controlled rent, which is approximately 2.9\% lower than benchmark's average free-market rent and about 9.1\% lower than the initial rent paid by new tenants in the controlled market in the benchmark.

For each year a tenant remains in the same unit, their rent decreases by $\kappa_{1}=\kappa_{2}$, mirroring the rate at which their unit deteriorates due to landlord neglect. Consequently, although their nominal rent declines during their tenancy, which results in reduced housing expense, their real rent remains constant.

Conversely, households drawing $\eta = 0$ are ineligible to rent in the controlled market but can opt for the free market. Additionally, regardless of the $\eta$ realization, households have the choice to purchase a house, reside with parents if $j<J_{p}$, or face homelessness.

To maintain the same proportion of the rental market under control as in the benchmark, we choose the probability of drawing $\eta = 1$, denoted as $\pi$. Beyond these adjustments, the model remains consistent with the baseline model's other assumptions.

\subsubsection{Aggregate results}

We present the aggregate findings of our initial experiment in the second column from the left in Table \ref{tab:aggregate_effects}. Given that we mechanically set the average controlled rent in the unit control experiment to match the average controlled rent in the benchmark, there would be no change in the average controlled rent. However, quantitatively, under unit control fewer controlled units would be allocated to low-income households, resulting in a worse misallocation of controlled units. This is because lotteries would exclude low-income households from the controlled market more substantially than the premium would exclude them under tenancy control. We illustrate this in panel (c) of Table \ref{tab:tenure_by_income_quintile}, which indicates a diminished proportion of low-income households renting in the controlled market compared to the benchmark. Thus, the homelessness rate would increase by 16.3\%.

Aggregate rental demand comprises both controlled-market demand and free-market demand. The demand within the controlled market is a function of the initial rent paid by new tenants upon entering controlled units. Conversely, free-market demand is a function of the average free-market rent. Consequently, aggregate rental demand is a three-dimensional object, where the aggregate rental quantity demanded is a function of the initial controlled rent and the average free-market rent.

Meanwhile, the aggregate quantity of rental housing supplied is a function of the marginal rent, as outlined in equation \eqref{eq:rental_supply}. Specifically, when the free market operates, the marginal rent aligns with the average free-market rent. Conversely, in the absence of a free market, the marginal rent corresponds to the average controlled rent.

To visually depict aggregate rental demand and supply, we use Figure \ref{fig:diagram_unit_control}. On the horizontal x-axis, we represent the quantity of rental housing ($X$), while the vertical y-axis represents average rents, where the average free-market rent determines both free-market demand and aggregate rental supply. To render aggregate rental demand in a comprehensible two-dimensional form, we plot it conditional on some initial controlled rent.

On the demand side, the solid red line, denoted as $D_{1}(p_{r,0}=p^{1}_{r,0})$, represents the benchmark's aggregate rental demand for an initial controlled rent, $p_{r,0}$, that is equal to the benchmark equilibrium initial controlled rent, $p^{1}_{r,0}$. The immediate rent discounts in the controlled market resulting from unit control, which take nearly 11 years to materialize in the benchmark with tenancy control, would spur an increase in aggregate rental demand at the intensive margin among controlled renters. This would consequently shift aggregate rental demand from $D_{1}(p_{r,0}=p^{1}_{r,0})$ to the dotted red line, $D_{2}(p_{r,0}=\overline{p}^{1}_{r,0}=\overline{p}^{2}_{r,0})$. $\overline{p}^{2}_{r,0}$ denotes the average controlled rent in the equilibrium under tenancy control, which is equal to the average controlled rent in the benchmark, $\overline{p}^{1}_{r,0}$. $D_{2}(p_{r,0}=\overline{p}^{1}_{r,0})$ represents aggregate rental demand under unit control when benchmark renters are not permitted to transition to homeownership, thus capturing the intensive margin. 

On the other hand, because unit control would not offer incumbent renters further discounts as they remain longer in their units, unit control would discourage long stays in the same unit relative to tenancy control. This would reduce aggregate rental demand at the extensive margin, shifting it from the dotted line $D_{2}(p_{r,0}=\overline{p}^{1}_{r,0})$, to the dashed line $D_{3}(p_{r,0}=\overline{p}^{1}_{r,0})$, which is the final location of aggregate rental demand.  

Overall, the impact coming from the intensive margin would outweigh the impact arising from extensive margin, causing aggregate rental demand to rise, resulting in average free-market rents rising by 2.4\%, from $\overline{p}^{1}_{r,1}$ to $\overline{p}^{2}_{r,1}$. This increase in free-market rents reduce the relative cost of homeownership for high-income households renting in the free market. This, combined with the reduced incentive to stay longer in the same unit, would lead more high-income households to opt for homeownership, resulting in a 0.15 percentage point increase in the homeownership rate. The increase in homeownership rates, combined with the overall increase in aggregate rental demand, would push house prices up by 0.4\%. Lastly, because under unit control, the average free-market rent would climb to $\overline{p}^{2}_{r,1}$ and the average controlled market rent would remain consistent with the benchmark, average rents under unit control would rise by 0.74\%.

On the rental supply side, the increase in house prices and property taxes\footnote{Property taxes would rise even though house prices would increase because homelessness would rise enough to offset the higher house prices.} would trigger a leftward shift in aggregate rental supply from $S_{1}$ to $S_{2}$, as outlined in equation \eqref{eq:rental_supply}. Conversely, akin to the findings in \citet{fallis-1985}, changes in the marginal rent would prompt \textsl{movements along the aggregate rental supply curve}. Overall, the higher average free-market rents would outweigh the leftward shift in supply, resulting in an expansion in the equilibrium quantity of rental housing, from $X^{1}$ to $X^{2}$, and a higher quantity of controlled rentals, moving from the benchmark quantity of, $\underline{X}^{1}_{0}$, to the equilibrium quantity, $\underline{X}^{2}_{0}$, under unit control. 

\subsubsection{Welfare results}

From the point of view of welfare, despite the fewer controlled units allocated to low-income households and the 2.4\% increase in average free-market rents, switching from tenancy control to unit control would only marginally decrease newborns' welfare by 0.06\%. However, these losses would be widespread, with only less than 1\% of newborns experiencing welfare gains. Welfare changes would be modest, primarily because unit control would promptly offer substantial rent discounts to renters fortunate enough to have access to the controlled market.

To understand the welfare implications of general equilibrium effects, we conduct a partial equilibrium simulation in which we maintain free rental prices at their benchmark levels but adjust other prices to align with long-run equilibrium under unit control. This approach allows us to isolate the impact of changes in average free-market rents on welfare outcomes. We find that when average free-market rents are held constant at their benchmark levels in this partial equilibrium scenario, welfare actually slightly increases, and three quarters of newborns see an improvement in their welfare. This indicates that the welfare gains due to the immediate benefits of reduced rents under unit control would outweigh the welfare losses resulting from the less desirable allocation of controlled units. Moreover, it also implies that it is critical to consider the general equilibrium effects of unit control on rental demand and average free-market rents. These general equilibrium impacts are significant enough to counterbalance the gains we observe in the partial equilibrium analysis.

We now turn to discussing distributional outcomes. In the second row of Table \ref{tab:distributional_impact}, we provide a detailed breakdown of welfare changes for newborns across different income brackets. We segment the income distribution into five quintiles. As the table demonstrates, welfare losses would be evenly distributed across all income quintiles. Low-income households would predominantly lose due to the allocation of fewer controlled units to them under unit control. Additionally, unlucky low-income households would encounter higher and volatile rents in the free market, along with an increased risk of homelessness. These factors would overshadow the benefits stemming from immediate rent discounts accessible to fortunate low-income households in the controlled market. 

Conversely, high-income households would primarily experience losses due to higher house prices and increased average free-market rents. Given their tendency to swiftly transition to homeownership, the advantages of immediate rent discounts and a greater allocation of controlled units to them would be outweighed by the aforementioned drawbacks. Mid-income households would face a combination of these challenges, resulting in overall welfare losses.

\subsubsection{Results summary}

To sum up, moving from tenancy control to unit control would lead to higher average free-market rents and an expansion in the equilibrium aggregate rental quantity. It would raise house prices, homeownership, and homelessness, as well as result in a slight decline in welfare, which would be evenly distributed across all income quintiles. 

\subsection{Benchmark vs. full decontrol}

\subsubsection{Aggregate results}

The third column of Table \ref{tab:aggregate_effects} delineates changes in aggregate variables in the long-run equilibrium under full decontrol compared to the benchmark, and Figure \ref{fig:diagram_full_decontrol} graphically shows the impact of moving from the benchmark to full decontrol. As in the previous experiment, we collapse aggregate rental demand to two dimensions by plotting it conditional on some initial controlled rent.

On the side of demand, mirroring the previous experiment, the solid red line, $D_{1}(p_{r,0}=p^{1}{r,0})$, illustrates the benchmark's aggregate rental demand given the benchmark equilibrium initial controlled rent, $p^{1}_{r,0}$. Full decontrol would elicit two effects on rental demand. On the one hand, the removal of controls in the rental market would drive down initial rents for renters, as average free-market rents are lower than initial controlled rents in the benchmark. This would stimulate aggregate rental demand at the intensive margin, causing it to shift from $D_{1}(p_{r,0}=p^{1}_{r,0})$ to the dotted red line $D_{2}(p_{r,0}=\infty)$. Here, $p_{r,0}=\infty$ signifies that the cost to rent in the controlled market is infinite, effectively indicating the shutting off of the controlled market. $D_{2}(p_{r,0}=\infty)$ represents aggregate rental demand under full decontrol assuming that no renters from the benchmark are permitted to transition to homeownership, thus capturing the intensive margin.

On the other hand, full decontrol would diminish renters' incentive to remain in the same rental unit for an extended period, prompting more renters to transition to homeownership. This extensive margin effect would consequently elevate homeownership rates, aggregate housing demand, and house prices. Additionally, it would induce a leftward shift in aggregate rental demand from $D_{2}(p_{r,0}=\infty)$ to the dashed red line, representing the final location of aggregate rental demand, $D_{3}(p_{r,0}=\infty)$. Here, $D_{3}(p_{r,0}=\infty)$ characterizes aggregate rental demand under full decontrol when benchmark renters can switch to homeownership.

In the grand scheme, the intensive margin effect would surpass the extensive margin effect, leading to a boost in aggregate rental demand among existing renters. This would result in a 0.1\% uptick in the average free-market rent, moving it from $\overline{p}^{1}_{r,1}$ to $\overline{p}^{2}_{r,1}$. Moreover, it would lead to a 2.1\% increase in the overall average rent, moving it from $\overline{p}^{1}_{r}$ to $\overline{p}^{2}_{r,1}$, because under full decontrol, the average free-market rent would become the average rent.

Moreover, in the absence of the controlled market option in our full decontrol equilibrium, more households would be exposed to rental price volatility and higher average rents. This exposure would culminate in a substantial 32.2\% surge in the homelessness rate.

We now turn to the dynamics of the aggregate rental supply. Since both house prices and property tax rates\footnote{Property tax rates would escalate despite the higher house prices because of the significant increase in homelessness, which would outweigh the increase in house prices.} would rise, aggregate rental supply would shift leftward from $S_{1}$ to $S_{2}$. However, the slight increase in average free-market rents would precisely counterbalance the leftward shift in supply. As a result, there would be no change in the equilibrium aggregate quantity of rental housing supplied, maintaining $X^{1}=X^{2}$.

\subsubsection{Welfare results}

Moving on to welfare, we find that decontrolling would adversely impact all newborns, reducing their welfare by 0.31\%. Additionally, it would lead to a decrease in welfare across the entire population by 0.14\%, with only 35.2\% of the population benefiting. These welfare losses would stem from two primary factors. Firstly, the removal of the controlled market would eliminate a valuable option for low-income households. Secondly, it would result in increased house prices and average free-market rents, adversely affecting both prospective homeowners and renters who previously rented in the free market.

Turning to distributional outcomes, in the third row of Table \ref{tab:distributional_impact}, we detail welfare changes across income quintiles. Under full decontrol, the reduction in newborns' welfare would be fairly uniform across the income distribution. Low-income households would primarily lose due to the elimination of the controlled market option, which is advantageous for them because they tend to rent in the controlled market in the benchmark. Conversely, high-income households would experience losses primarily because of the higher average free-market rents, since they tend to rent in the free market, and because of the more expensive house prices, rendering homeownership less accessible to them. Mid-income households would face a combination of these factors, contributing to their welfare losses.

\subsubsection{Results summary}

In summary, full decontrol would lead to an increase in average rents and average free-market rents while maintaining equilibrium aggregate rental quantity. This would result in higher house prices and greater rates of homeownership homelessness. Additionally, it would lead to a welfare loss, which would be rather evenly distributed across households in different income quintile. 

\subsection{Benchmark vs. full tenancy control}\label{subsection:full_control}

\subsubsection{Aggregate results}

In our third experiment, we analyze the long-run equilibrium under full tenancy control in comparison to the benchmark. The fourth column from the left in Table \ref{tab:aggregate_effects} delineates the changes in aggregate variables, while Figure \ref{fig:diagram_full_control} presents a qualitative diagram elucidating the ramifications of moving from the benchmark to full tenancy control.

The implementation of full tenancy control would result in all renters renting in the controlled market. Consequently, households would no longer be susceptible to the volatility of free-market rents, leading to a 2\% decrease in the homelessness rate.

As observed in the preceding experiments, the solid red line in Figure \ref{fig:diagram_full_control}, denoted as $D_{1}(p_{r,0}=p^{1}_{r,0})$, represents aggregate rental demand conditional on the benchmark equilibrium's initial controlled rent, $p^{1}_{r,0}$. However, portraying aggregate rental demand under full tenancy control presents a challenge, because aggregate rental demand becomes the same as aggregate controlled rental demand, which exists only on the z-x plane, where the z-axis represents the initial controlled rent and the x-axis mirrors that of Figure \ref{fig:diagram_full_control}. To graphically depict equilibrium changes in aggregate rental demand under full tenancy control, we resort to using a point on the graph.

Starting with changes in rental demand, under full tenancy control, renters who previously opted for the free market would now encounter higher initial rents in the controlled market. This would result in a reduction in their rental demand at the intensive margin. We illustrate this change using point $D_{2}(\overline{p}_{r,1}=\infty,p_{r,0}=\hat{p}^{2}_{r,0})$, which signifies the intermediate position of aggregate rental demand under the premise that benchmark renters are prohibited from transitioning to homeownership, and vice versa, effectively capturing the intensive margin. The term $\overline{p}_{r,1}=\infty$ encapsulates the concept that in this experiment, the free rental market ceases to exist, which is akin to having an infinitely high cost associated with renting in that market. Meanwhile, $\hat{p}^{2}_{r,0}$ denotes the equilibrium initial controlled rent under full tenancy control, when no benchmark renters are permitted to transition to homeownership and no benchmark homeowners are allowed to become renters. Point $D_{2}(\overline{p}_{r,1}=\infty,p_{r,0}=\hat{p}^{2}_{r,0})$ says that for the pair ($\hat{\bar{p}}^{2}_{r,1},\hat{p}^{2}{r,0})$, the quantity demanded equals the quantity supplied. Here, $\hat{\bar{p}}^{2}_{r,1}$ and $\hat{p}^{2}_{r,0}$ represent the equilibrium average controlled rent and initial controlled rent under full tenancy control, respectively, in a scenario where benchmark renters cannot transition to homeownership, and contrariwise.

On the other hand, full tenancy control would entice more high-income households to remain longer in the rental market compared to the benchmark. This would lead to a 0.26 percentage point reduction in the homeownership rate. Consequently, aggregate rental demand would shift rightward from point $D_{2}(\overline{p}_{r,1}=\infty,p_{r,0}=\hat{p}^{2}_{r,0})$ to point $D_{3}(\overline{p}_{r,1}=\infty,p_{r,0}=p^{2}_{r,0})$, due to this extensive margin. The latter point represents the final position of aggregate rental demand when benchmark renters are permitted to transition to homeownership, and vice versa, effectively capturing the extensive margin. The point $D_{3}(\overline{p}_{r,1}=\infty,p_{r,0}=p^{2}_{r,0})$ means that for the rent bundle ($\overline{p}^{2}_{r,0},p^{2}_{r,0})$, the quantity demanded equals the quantity supplied, where $\overline{p}^{2}_{r,0}$ and $p^{2}_{r,0}$ represent equilibrium average controlled rent and average initial rent, respectively, under full tenancy control.

Overall, the reduction in aggregate rental demand originating from the intensive margin would surpass the rise in aggregate rental demand from the extensive margin, resulting in a decrease in aggregate rental demand relative to the benchmark. Since both aggregate rental demand and homeownership rates would decrease, house prices would also experience a decline of 0.3\%.

Moving on to rental supply, in contrast to the unit control and full decontrol experiments, the decline in house prices alongside unchanged property taxes, attributed to the marginal decrease in the homelessness rate, would lead to a rightward shift in aggregate rental supply from $S_{1}$ to $S_{2}$. Recall that, under full tenancy control, the average controlled rent would become the marginal rent for the rental company. Since the average controlled rent would be lower than the benchmark average free-market rent and would outweigh the rightward shift in rental supply, the equilibrium quantity of rental housing would decrease by 0.6\% from $X^{1}$ to $X^{2}$.

Because the average controlled rent would rise by 0.7\% from $\overline{p}^{1}_{r,0}$ to $\overline{p}^{2}_{r,0}$, there would be an increase in the equilibrium quantity of controlled rentals from $\underline{X}^{1}_{0}$ to $X^{2}$. The rise in the average controlled rent can be attributed to several factors. In the benchmark, high-income households often transition to homeownership more swiftly than low-income households, resulting in their predominant presence in the free rental market. However, when the free market option is removed, more high-income households would be compelled to rent in the controlled market. Consequently, this shift would shorten the average duration of rental tenure in the controlled market, leading to an increase in both the average controlled rents and the quantity of controlled rentals. Thus, despite the surge in demand within the controlled market, the increase in the quantity supplied would perfectly offset the additional demand. This means that the initial controlled rents would not need to change at all in order to equate demand and supply within the controlled market. 

\subsubsection{Welfare results}

In terms of welfare, we find a slight decline in the welfare of newborns, equal to 0.02\%, although the average welfare across all households would slightly increase, by 0.16\%. The decrease in newborns' welfare would primarily stem from the removal of the free market option, which would reduce available rental choices for households. The presence of a free market option benefits mainly higher-income households, but occasionally some low-income households also prefer renting in the free market to avoid premium costs because some very low-income households cannot afford to pay the premium to rent in the controlled market or have to spend almost their entire income to access the controlled market. Consequently, less than a quarter of all newborns would benefit from this policy change.

Conversely, the average welfare across all households would slightly rise because our calculation of average welfare includes households that have already been in the controlled market and have accumulated some longer-term benefits of rent control. However, these benefits would extend only to approximately 27\% of the population.

Lastly, in the fourth row of Table \ref{tab:distributional_impact}, we present a breakdown of the welfare gains and losses of newborns based on their position in the income distribution. Welfare losses are evident across all income quintiles. This would be attributed to the removal of the free market, a rental option that proves beneficial for numerous households, as we detailed above.

\subsubsection{Results summary}

In conclusion, full tenancy control would result in a reduction of both average rents and equilibrium aggregate rental quantity. This would lead to a slight decrease in house prices, and lower rates of homeownership and homelessness. Additionally, it would result in a slight decline in welfare, mostly impacting low- and high-income households. 

\subsection{Optimal degree of tenancy control}\label{subsection:optimal_control}

We have demonstrated that, in our model, both full decontrol and full tenancy control would result in overall welfare losses. This prompts the next logical inquiry: what level of tenancy control would maximize welfare? To answer this, we conduct a comprehensive grid search. We systematically vary the share of the market subject to tenancy control, and assess the welfare corresponding to the long-run equilibrium of each case.

In our model, the welfare-maximizing degree of tenancy control is 65\%, which closely aligns with our benchmark in which approximately 70\% of the rental market falls under tenancy control. This level of control would yield only a marginal 0.01\% increase in welfare compared to the benchmark. The optimal degree of tenancy control strikes a delicate balance between two competing objectives.

The first objective is to maintain a controlled market with affordable premiums, primarily benefiting low-income households by easing their housing costs and facilitating access to this segment of the market. When the government decreases the proportion of the controlled market, competition for controlled units intensifies, resulting in higher premiums that burden low-income households.

The second aim is to keep rents in the free market affordable for households who prefer renting outside the controlled market, which predominantly includes higher-income households, but also some low-income households who opt for this option as we explained in Section \ref{subsection:full_control}. Achieving this goal involves reducing the size of the controlled market. However, beyond a certain threshold, an expansion of the controlled market requires average free-market rents to rise in order to encourage more households to rent in the controlled market, which hurts those households who would otherwise prefer to rent in the free market. In our model, we find that these trade-offs are effectively balanced when 65\% of the rental market is subject to tenancy control.

For completeness, we include the aggregate results of this last experiment in the fifth column of Table \ref{tab:aggregate_effects} and the distributional outcomes in the fifth row of Table \ref{tab:distributional_impact}. However, since outcomes in this welfare-maximizing long-run equilibrium are negligibly different from the benchmark, for brevity's sake we skip further discussing them.

Given the proximity of the welfare-maximizing degree of tenancy control to the benchmark level, which we have already compared to unit control, we also opt not to delve further into its differences relative to unit control. However, we did replicate the exercise with unit control and observed outcomes akin to those of our initial experiment.

\section{Conclusion}\label{sec:conclusion}

We constructed a general equilibrium model encompassing both tenancy control and unit control --- two prevalent rent-control policies --- to analyze their respective impacts on welfare and equilibrium prices. Leveraging household-level 2016 Census data and household-level Survey of Financial Security (SFS) data from Toronto, where approximately 71\% of rental units fall under tenancy control while the remainder remain uncontrolled, we calibrated our model to reflect the local housing market dynamics.

Utilizing the calibrated model, we conducted four distinct experiments to evaluate the impacts of various rent-control policies. In the first experiment, we compared the benchmark against the long-run equilibrium under unit control, maintaining an identical degree of control and average controlled rents. The second experiment compared the long-run equilibrium under full rent decontrol with the benchmark. In the third experiment, we contrasted the benchmark with the long-run equilibrium under full tenancy control. Our final experiment involved pinpointing the welfare-maximizing degree of tenancy control in the long run.

Across these experiments, we obtain several main findings. Unit control would result in fewer controlled rental units being allocated to low-income households compared to tenancy control. This is because under unit control, lotteries determine which households can rent in the controlled market. In contrast, under tenancy control, despite the lower bound on rental housing consumption, the price mechanism is still able to allocate more controlled units to low-income households with higher willingness to pay. Furthermore, unit control would lead to an increase in aggregate rental demand, consequently elevating average free-market rents, and thus marginally diminishing welfare and exacerbating homelessness relative to the benchmark. Both full decontrol and full control would result in a reduction in welfare across the entire income spectrum relative to the benchmark, primarily due to the elimination of viable rental options for households. Finally, the welfare-maximizing degree of tenancy control aims to strike a balance between reducing average free-market rents and maintaining affordable initial rents in the controlled market. In our model, we identified that under the welfare-maximizing long-run equilibrium 65\% of the rental market would be under tenancy control, a degree closely aligned with our benchmark.

Our analysis underscores critical impacts of rent-control policies on housing access for low-income households and homelessness. Policymakers should carefully weigh how different policy approaches may influence the ability of vulnerable populations to secure affordable housing, especially in high-demand urban areas grappling with pressing housing affordability challenges. Part of this requires a firm grasp of the trade-offs between immediate rent discounts offered under unit control that come at the potential exclusion of more low-income households from accessing controlled units due to shortages, which raises homelessness relative to tenancy control.

Empirical evidence from studies such as \citet{nagy-1997}, \citet{early-2000}, and \citet{cjq-2023} in New York City, where unit control is the prevailing rent-control policy, suggests that new tenants of controlled units often face premiums relative to the free market. This implies widespread non-compliance among landlords whose units are under control, inadvertently nudging the equilibrium closer to the tenancy control equilibrium. Although we do not condone breaking the law, given that tenancy control would improve welfare outcomes, this non-compliance may actually be to the overall benefit of society.

Lastly, we see an opportunity for future research to concentrate on the impacts of tenancy control on evictions relative to unit control. Because tenants under tenancy control receive larger rent discounts the longer they stay in the unit, landlords have greater incentives to evict longer-standing tenants than landlords who operate under unit control. It would be important for policymakers to consider the magnitude of the impact of this channel on welfare outcomes as well.

\bibliographystyle{aer}
\bibliography{reference_files/bib}

\vspace{1cm}

\begin{table}[H]
    \centering
    \caption{Externally assigned parameter values}
    \vspace{-1.25cm}
    \label{tab:calibration_external}
    \input{tables/calibration_external}
    \vspace{-0.5cm}
    \begin{fignote2}{1.11\textwidth}
    \textit{Notes:} Table \ref{tab:calibration_external} presents the values of the externally calibrated parameters. Part (i) delineates the external calibration for non-housing parameters, while part (ii) outlines the external calibration of housing parameters. Additionally, part (iii) provides the external calibration of mortgage-related parameters.
    \end{fignote2}
\end{table}

\begin{table}[H]
    \centering
    \caption{Internally assigned parameter values}
    \vspace{-1.25cm}
    \label{tab:calibration_internal}
    \input{tables/calibration_internal}
    \vspace{-0.5cm}
    \begin{fignote2}{1.08\textwidth}
    \textit{Notes:} Table \ref{tab:calibration_internal} displays the values of the internally calibrated parameters. Our benchmark is closely aligned with all of our targets. From the 2019 SFS, we compute the percentage of homeowners with mortgages and the aggregate bequests to aggregate net wealth ratio. From the 2016 Census, we obtain data for Toronto CMA regarding the total rent to total income ratio, the average income tax rate, the labor income Gini coefficient, the percentage of renters in the controlled market, the percentage of renters spending over 50\% of their income on rent, the homeownership rate, and the percentage of 26-34-year-old households living with their parents. Additionally, using 2016 Census and 2016 CREA data for Toronto CMA, we calculate the average free-market rent to housing price ratio. For the tax incidence of renters, we use estimates from \citet{cy-1994}, \citet{susin-2002}, \citet{fack-2006}, and \citet{sy-2020}. Lastly, we obtain the homelessness rate from the ``Homelessness Hub".
    \end{fignote2}
\end{table}

\begin{table}[H]
    \centering
    \caption{Aggregate non-targeted statistics}
    \vspace{-0.5cm}
    \label{tab:aggregate_non_targreted_statistics}  \input{tables/aggregate_non_targeted_statistics.tex}
    \vspace{-0.45cm}
    \begin{fignote2}{0.9\textwidth}
      \textit{Notes:} Table \ref{tab:aggregate_non_targreted_statistics} presents the model's performance on non-targeted aggregate statistics compared to the data. The second column from the left displays the model's performance, the third column includes estimates from the literature or the data, and the last column provides the source. ``--" indicates that the statistic is unavailable.
    \end{fignote2}
\end{table}

\begin{table}[H]
    \centering
    \caption{Rental market tenure by income quintile (all households)}
    \vspace{-0.5cm}
    \label{tab:tenure_by_income_quintile}    
    \input{tables/tenure_by_income_quintile_everyone.tex}
    \vspace{-0.45cm}
    \begin{fignote2}{0.66\textwidth}
      \textit{Notes:} Table \ref{tab:tenure_by_income_quintile} illustrates rental market tenure by income quintile for all households. Panel (a) presents these statistics in the data, which we computed using 2016 Census data, panel (b) showcases the benchmark's performance, and panel (c) displays outcomes of our unit control exercise. The second column from the left represents the lowest income quintile, while the first column from the right signifies the highest quintile. The first row in each panel indicates the percentage of households in each quintile renting in the controlled market. For instance, in panel (a), 28.88\% of all households in the second quintile rent in the controlled market. The second row in each panel indicates the percentage of households in each quintile renting in the free market.  
    \end{fignote2}
\end{table}

\begin{table}[H]
    \centering
    \caption{Aggregate effects generated by changes in rent-control policy}
    \vspace{-0.5cm}
    \label{tab:aggregate_effects}
    \input{tables/aggregate_effects.tex}
    \vspace{-0.45cm}
    \begin{fignote2}{1.02\textwidth}
      \textit{Notes:} In Table \ref{tab:aggregate_effects}, the second column contrasts the benchmark with the unit control long-run equilibrium. The third column juxtaposes the benchmark with the long-run equilibrium under full decontrol. The fourth column compares the benchmark with the long-run equilibrium under full tenancy control. The last column compares the benchmark with the welfare-maximizing degree of tenancy control in the long run. ``—-" indicates incompatible statistics. We assess welfare in non-housing consumption-equivalent terms.
    \end{fignote2}
\end{table}

\begin{table}[H]
    \centering
    \caption{Distributional impacts of rent-control policy changes on newborns' welfare, broken down by income quintile ($\%\Delta$ in welfare)}
    \vspace{-0.5cm}
    \label{tab:distributional_impact}
    \input{tables/distributional_impacts.tex}
    \vspace{-0.45cm}
    \begin{fignote2}{0.77\textwidth}
      \textit{Notes:} Table \ref{tab:distributional_impact} presents welfare gains and losses for newborns, comparing unit control, full decontrol, and full tenancy control to the benchmark. We measure welfare in non-housing consumption-equivalent terms.
    \end{fignote2}
\end{table}

\begin{figure}[H]
  \centering
    \caption{Diagram illustrating rent trajectories under tenancy control based on rental market selection and duration in the same rental unit}
    \vspace{-0.37cm}
    \label{fig:diagram_rent_path}
    \includegraphics[width=\textwidth]{figures/path_of_rent_diagram.pdf}
    \vspace{-1.3cm}
    \begin{fignote2}{1.0\textwidth}
      \textit{Notes:} The downward-sloping blue line depicts rental prices paid by tenants residing in controlled units as a function of their duration of occupancy in those units, where the initial rent in the controlled market is $p_{r,0}$ and $\overline{p}_{r,0}$ is the average rent in that market. Conversely, the horizontal red line represents the average rental price faced by tenants opting for the free market, denoted as $\overline{p}_{r,1}$. The dashed red line represents free-market rent fluctuations tenants in the free market may face. Notably, the average free-market rent remains constant regardless of tenancy duration. For clarity in this illustration, we abstract from the deterioration of units in the controlled market. Early during a tenancy, tenants in the free market pay lower average rents than those in controlled units, but later in the tenancy this reverses. Why is this the case? If initial rents in the controlled market were equal to or lower than average free-market rents, then all renters would logically opt for controlled rentals. To entice some tenants towards the free market, those in controlled units must initially pay a premium to enjoy the safeguarding against rent fluctuations and lower long-term rental prices. Tenants anticipating short-term stays and willing to tolerate risk, opt for the free market to avoid this premium. Conversely, tenants planning extended stays and having low risk tolerance, are willing to pay the initial premium to capitalize on the advantages of rent control.      
    \end{fignote2}
\end{figure}

\begin{figure}[H]
  \centering
    \caption{Lifecycle statistics}
    \vspace{-0.2cm}
    \label{fig:life_cycle}
    \includegraphics[width=\textwidth]{figures/lifecycle.pdf}
    \vspace{-0.45cm}
    \begin{fignote2}{1\textwidth}
      \textit{Notes:} In Figure \ref{fig:life_cycle}, the two panels illustrate the model's performance on lifecycle statistics compared to the data. Panel (a) displays the average homeownership rate for each age cohort, while panel (b) presents the average loan-to-value (LTV) ratios for homeowners across 12 age groups, as the 2019 SFS provides household maintainer's age in 5-year ranges. Model outputs are depicted in orange, while data points are represented in blue. We source homeownership rates by age for Toronto CMA from the 2016 Census, and obtain data on the average LTV ratios from the 2019 SFS.
    \end{fignote2}
\end{figure}

\begin{figure}[H]
  \centering
    \caption{Rent-to-income ratio cumulative distribution function}
    \vspace{-0.2cm}
    \label{fig:rent_to_income_cdf}
    \includegraphics[width=\textwidth]{figures/rent_income_cdf.pdf}
    \vspace{-0.45cm}
    \begin{fignote2}{1\textwidth}
      \textit{Notes:} In Figure \ref{fig:rent_to_income_cdf}, we depict the cumulative distribution function (CDF) of the rent-to-income ratio for renters. We represent the model's output in orange, and the data in blue. We calculate the rent-to-income CDF using data from the 2016 Census for Toronto CMA.
    \end{fignote2}
\end{figure}

\begin{figure}[H]
  \centering
    \caption{Qualitative diagram illustrating the impact of switching from the benchmark with tenancy control to unit control}
    \vspace{-0.36cm}
    \label{fig:diagram_unit_control}
    \includegraphics[width=\textwidth]{figures/unit_control_diagram.pdf}
    \vspace{-0.6cm}
    \begin{fignote2}{1.0\textwidth}
      \textit{Notes:} In Figure \ref{fig:diagram_unit_control}, we present a qualitative diagram illustrating the impact of switching from the benchmark to unit control. The x-axis represents the aggregate quantity of rental housing, denoted as $X$. The y-axis represents average rents; the average free-market rent determines the quantity demanded in the free-market and the aggregate quantity supplied. The benchmark aggregate rental supply is the solid blue line, denoted as $S_{1}$. The aggregate rental supply under unit control is the dashed blue line, denoted as $S_{2}$. The benchmark aggregate rental demand is the solid red line, $D_{1}(p_{r,0}=p^{1}_{r,0})$, where $p^{1}_{r,0}$ is the benchmark equilibrium initial controlled rent. The final aggregate rental demand under unit control is the dashed red line, denoted as $D_{3}(p_{r,0}=\overline{p}^{1}_{r,0})$, where $\overline{p}^{1}_{r,0}$ is the equilibrium average controlled rent under unit control, which is also the initial rent in the controlled market under unit control. The dotted red line, $D_{2}(p_{r,0}=\overline{p}^{1}_{r,0})$, represents intermediate aggregate rental demand under the unit control equilibrium. This curve measures aggregate rental demand when benchmark renters are not allowed to switch to homeownership and homeowners are not allowed to become renters, thus capturing the intensive margin. In the benchmark equilibrium, $E_{1}$, we denote the controlled rental quantity as $\underline{X}^{1}_{0}$, aggregate rental quantity as $X^{1}$, average free-market rent as $\overline{p}^{1}_{r,1}$, average controlled rent as $\overline{p}^{1}_{r,0}$, and average rent as $\overline{p}^{1}_{r}$. In the unit control equilibrium, $E_{2}$, we denote controlled rental housing quantity as $\underline{X}^{2}_{0}$, aggregate rental quantity as $X^{2}$, average controlled rent as $\overline{p}^{2}_{r,0}=\overline{p}^{1}_{r,0}$, and average free-market rent as $\overline{p}^{2}_{r,1}$.
    \end{fignote2}
\end{figure}

\begin{figure}[H]
  \centering
    \caption{Qualitative diagram illustrating the impact of full rent decontrol}
    \vspace{-0.36cm}
    \label{fig:diagram_full_decontrol}
    \includegraphics[width=\textwidth]{figures/full_decontrol_diagram.pdf}
    \vspace{-0.6cm}
    \begin{fignote2}{1.0\textwidth}
      \textit{Notes:} In Figure \ref{fig:diagram_full_decontrol}, we present a qualitative diagram illustrating the impact of full decontrol compared to the benchmark. The x- and y-axes are the same as in Figure \ref{fig:diagram_unit_control}. The benchmark aggregate rental supply is the solid blue line, denoted as $S_{1}$. The aggregate rental supply under full rent decontrol is the dashed blue line, which we denote as $S_{2}$. The benchmark aggregate rental demand is the solid red line, $D_{1}(p_{r,0}=p^{1}_{r,0})$, where $p^{1}_{r,0}$ is the benchmark equilibrium initial controlled rent. The final aggregate rental demand under full rent decontrol is the dashed red line, denoted as $D_{3}(p_{r,0}=\infty)$, where $p_{r,0}=\infty$ represents the idea that the controlled rental market has been shut off, which is equivalent to saying that the initial controlled rent is infinity. The dotted red line, $D_{2}(p_{r,0}=\infty)$, represents intermediate aggregate rental demand under the full decontrol equilibrium, which considers what aggregate rental demand would be if no benchmark renters were allowed to switch to homeownership and no homeowners were allowed to become renters. In the benchmark equilibrium, $E_{1}$, we denote the controlled rental quantity as $\underline{X}^{1}_{0}$, aggregate rental quantity as $X^{1}$, average free-market rent as $\overline{p}^{1}_{r,1}$, average controlled rent as $\overline{p}^{1}_{r,0}$, and average rent as $\overline{p}^{1}_{r}$. In the full decontrol equilibrium, $E_{2}$, we denote aggregate rental quantity as $X^{2}$, and average free-market rent as $\overline{p}^{2}_{r,1}$.
    \end{fignote2}
\end{figure}

\begin{figure}[H]
  \centering
    \caption{Qualitative diagram illustrating the impact of full tenancy control}
    \vspace{-0.36cm}
    \label{fig:diagram_full_control}
    \includegraphics[width=\textwidth]{figures/full_tenancy_control_diagram.pdf}
    \vspace{-0.8cm}
    \begin{fignote2}{1.0\textwidth}
      \textit{Notes:} Figure \ref{fig:diagram_full_control} provides a qualitative diagram depicting the impact of full tenancy control relative to the benchmark. The x- and y-axes are the same as in Figure \ref{fig:diagram_unit_control}. The benchmark aggregate rental supply is the solid blue line, denoted as $S_{1}$. The aggregate rental supply under full tenancy control is the dashed blue line, which we denote as $S_{2}$. As in the other figures, the benchmark aggregate rental demand is the solid red line, $D_{1}(p_{r,0}=p^{1}_{r,0})$, where $p^{1}_{r,0}$ is the benchmark equilibrium initial controlled rent. Once the full rental market is under control, however, aggregate rental demand disappears from the x-y plane. Why does it disappear? Under full tenancy control, aggregate rental demand is equal to aggregate controlled rental demand. Since aggregate controlled rental demand is a function of initial controlled rents, aggregate rental demand under full tenancy control is situated on a different plane than aggregate rental supply --- the z-x plane, where the z-axis represents initial controlled rents and the x-axis is the same as in this figure. In order to still graphically show it, we use a point on the x-y graph to represent aggregate rental demand. The final aggregate rental demand under full tenancy control is a point, which we denote as $D_{3}(\overline{p}_{r,1}=\infty,p_{r,0}=p^{2}_{r,0})$. $\overline{p}_{r,1}=\infty$ represents the idea that the free rental market has been banned, which is equivalent to saying that the cost to rent in that market is infinite. $p^{2}_{r,0}$ denotes the equilibrium initial controlled under full tenancy control. This point means that for the bundle ($\overline{p}^{2}_{r,0},p^{2}_{r,0})$, the quantity demanded equals to the quantity supplied, where $\overline{p}^{2}_{r,0}$ denotes the equilibrium average controlled rent under full tenancy control. Intermediate aggregate rental demand, denoted as point $D_{2}(\overline{p}_{r,1}=\infty,p_{r,0}=\hat{p}^{2}_{r,0})$, considers what aggregate rental demand would be if no benchmark renters were allowed to switch to homeownership and no benchmark homeowners were allowed to switch to renting. This point means that for the bundle ($\hat{\bar{p}}^{2}_{r,0},\hat{p}^{2}_{r,0})$, the quantity demanded equals to the quantity supplied, where $\hat{\bar{p}}^{2}_{r,0}$ and $\hat{p}^{2}_{r,0}$ denote what the equilibrium average controlled rent and initial controlled rent under full tenancy control would be if no benchmark renters were allowed to become homeowners and vice versa. In the benchmark equilibrium, $E_{1}$, we denote the controlled rental quantity as $\underline{X}^{1}_{0}$, aggregate rental quantity as $X^{1}$, average free-market rent as $\overline{p}^{1}_{r,1}$, average controlled rent as $\overline{p}^{1}_{r,0}$, and average rent as $\overline{p}^{1}_{r}$. In the full tenancy control equilibrium, $E_{2}$, we denote aggregate rental quantity as $X^{2}$, and average controlled rent as $\overline{p}^{2}_{r,0}$.  
    \end{fignote2}
\end{figure}
    
\newpage

\include{appendix}

\end{document}